\chapter{长上下文：让模型阅读整本书}
\label{ch:long-context}
%% ============================================================

早期的 Transformer 模型只能处理 512 或 2048 个 token，
大约一两页文字。这严重限制了模型的实用性：
它无法阅读一篇长论文，无法理解一个完整的代码仓库，
更无法在长篇对话中保持一致性。

2025--2026 年，长上下文能力取得了巨大突破。
下表展示了截至 2026 年初主要模型的上下文窗口对比：

\begin{center}\small
\renewcommand{\arraystretch}{1.3}
\begin{tabular}{lcll}
\toprule
\rowcolor{figLightGray}
\textbf{模型} & \textbf{上下文窗口} & \textbf{位置编码} & \textbf{发布时间} \\
\midrule
LLaMA 4 Scout     & 10M tokens  & iRoPE       & 2025.04 \\
Gemini 2.5 Pro    & 1M tokens   & 未公开       & 2025.03 \\
Gemini 3 Pro/Flash & 1--2M tokens & 未公开      & 2025--2026 \\
Claude Sonnet 4.5 / Opus 4.5 & 200K--1M tokens & 未公开 & 2025--2026 \\
GPT-5 系列        & 400K tokens & 未公开       & 2025--2026 \\
DeepSeek-V3.2     & 1M tokens   & YaRN + MLA   & 2026.02 \\
Kimi K2.5          & 256K tokens & RoPE 变体    & 2026.01 \\
Qwen3-Next        & 262K--1M tokens & YaRN 变体 & 2026 \\
GLM-4.7           & 200K tokens & RoPE 变体    & 2025.12 \\
Jamba 1.5-Large   & 256K tokens & RoPE + SSM   & 2024.08 \\
DeepSeek-V3       & 128K tokens & YaRN + MLA   & 2024.12 \\
LLaMA 3.1 405B    & 128K tokens & RoPE + 频率调整 & 2024.07 \\
Qwen 2.5          & 128K--1M tokens & YaRN    & 2024--2025 \\
\bottomrule
\end{tabular}
\end{center}

LLaMA 4 Scout 以 10M token（约 15,000 页）
创下了公开模型的最长上下文记录。
Google 的 Gemini 系列稳定支持 1--2M token。
Anthropic 的 Claude 系列在 2025--2026 年
将上下文扩展到了 200K--1M token。
OpenAI 的 GPT-5 系列达到了 400K token。

值得注意的是，\textbf{1M token 上下文在 2026 年初已成为事实上的标准}。
DeepSeek 在 2026 年 2 月 11 日悄然将 V3.2 的上下文窗口
从 128K 升级到 1M，没有发布公告，
用户在使用中自行发现了这一变化。
Kimi K2.5 支持 256K 原生上下文，
Qwen3-Next 原生支持 262K 并可扩展至 1M，
GLM-4.7 支持 200K 上下文和 128K 最大输出长度。
这意味着长上下文竞争的焦点
已经从“能否支持 1M”转向“能否有效利用 1M”。

一个值得关注的趋势是\textbf{输出长度}的增长。
传统上，长上下文主要指长\textbf{输入}，
模型读入大量文本并做理解和检索。
但 2026 年初，多个模型开始支持长\textbf{输出}，
GLM-4.7 支持 128K 的最大输出长度，
Claude Opus 4.5 支持 64K+ 的输出。
这使得模型不仅能“阅读整本书”，
还能“写出整本书”，
在保持前后一致性的前提下
生成数万 token 的连贯文本。
长输出对 KV cache 管理提出了不同于长输入的挑战：
在生成过程中，KV cache 持续增长，
内存压力随生成进度递增，
需要动态的内存管理策略来应对。

但这些数字背后隐藏着艰巨的工程和科学挑战，
让模型「看到」长文本和让模型「理解」长文本，
完全是两回事。
正如 Chroma 的 context rot 研究~\cite{chroma2025contextrot} 所揭示的：
\textbf{标称上下文长度和有效上下文利用率之间存在巨大鸿沟。}

\section{为什么长上下文很难？}

长上下文的困难不是单一的，而是四重挑战的叠加。

\subsection{计算复杂度：$O(n^2)$ 的诅咒}

注意力机制的计算复杂度是 $O(n^2)$：
序列长度翻倍，计算量翻四倍。
让我们把这个抽象的复杂度换算成具体的数字：

处理 4K token 时，每一层注意力需要计算
$4{,}096 \times 4{,}096 = 16.8M$ 个注意力分数对。
处理 128K token 时，这个数字变成
$131{,}072 \times 131{,}072 \approx 16.4B$，
\textbf{增长了近 1000 倍}。
如果模型有 80 层（如 LLaMA 70B），
那就是 $80 \times 16.4B \approx 1.3$ 万亿次注意力分数计算。

即使有 FlashAttention~\cite{dao2022flashattention} 这样的优化
（将注意力的 IO 复杂度从 $O(n^2)$ 降到 $O(n^2/M)$，
其中 $M$ 是 SRAM 大小），
计算量本身的 $O(n^2)$ 增长是无法改变的。
128K token 的计算量就是 4K token 的 $1024\times$，
这是数学决定的，没有捷径。

\subsection{内存瓶颈：KV Cache 的爆炸}

比计算更致命的是内存。
在推理时，Transformer 需要缓存所有已处理 token 的
Key 和 Value 向量（称为 KV cache），
以便后续 token 做注意力计算。

KV cache 的大小可以这样估算：
\begin{equation}
  \text{KV cache} = 2 \times n_{\text{layers}} \times n_{\text{kv\_heads}}
  \times d_{\text{head}} \times n_{\text{tokens}} \times \text{bytes per element}
  \label{eq:kv-cache-size}
\end{equation}
其中因子 2 来自 Key 和 Value 两组向量，
$n_{\text{kv\_heads}}$ 是 KV 头的数量（注意不是 query 头的数量：
在使用 GQA 的模型中两者不同）。

LLaMA 70B 使用 GQA（Grouped Query Attention），
有 64 个 \textbf{query 头}但只有 8 个 \textbf{KV 头}
（每 8 个 query 头共享一组 KV 投影）。
因此 KV cache 计算直接使用 $n_{\text{kv\_heads}} = 8$。

对于 LLaMA 70B（80 层，8 个 KV heads，head 维度 128）
处理 128K token，使用 BF16（2 bytes）：
\[
  2 \times 80 \times 8 \times 128 \times 131{,}072 \times 2
  \approx 21 \text{ GB}
\]
仅一个用户的一次 128K 请求，
光是 KV cache 就要占用约 21\,GB，
超过单块 H100（80\,GB）显存的四分之一，
模型权重还没算。
如果不使用 GQA（即 64 个 KV heads），
KV cache 将膨胀到约 171\,GB，远超单块 GPU 容量。

\begin{warnbox}[KV Cache 是长上下文推理的真正瓶颈]
很多人以为长上下文的瓶颈是计算。
实际上，对于推理来说，\textbf{内存才是硬约束}。
计算可以通过更多 GPU 和更长时间来解决，
但 KV cache 必须驻留在 GPU 显存中，
而单块 GPU 的显存是有限的。
这就是为什么 KV cache 的压缩和分布式管理
是长上下文推理工程中最核心的问题。
\end{warnbox}

KV cache 的爆炸对\textbf{并发处理能力}的影响更为深远。
在推理服务中，一块 GPU 通常需要同时处理多个用户请求。
如果每个 128K 请求需要 21\,GB KV cache，
一块 80\,GB 的 H100 在装载模型权重（如 30\,GB）后
只剩约 50\,GB 给 KV cache，
最多同时服务 2 个 128K 请求。
而如果每个请求只用 4K 上下文，
KV cache 仅约 0.7\,GB/请求，
同一块 GPU 可以同时服务 70+ 个请求。
\textbf{长上下文将并发能力降低了 35 倍}。

这就是为什么 KV cache 压缩技术如此关键。
DeepSeek-V3 使用的 MLA 将 KV cache 压缩了约 $8\times$
（通过潜在空间投影），
使得 128K 上下文的 KV cache 从 $\sim$21\,GB
降低到约 $\sim$2.6\,GB，
极大地提升了长上下文推理的并发处理能力。

除了 MLA，还有几种 KV cache 压缩策略被广泛使用：

\textbf{KV Cache 量化}。
将 KV cache 从 BF16（2 bytes）压缩到 INT8（1 byte）
或 INT4（0.5 bytes），直接减少 2--4 倍内存。
KV cache 的量化比模型权重量化更具挑战性，
因为注意力分数对 Key 向量的微小误差非常敏感，
一个 Key 向量的量化误差可能改变注意力的排序，
导致模型关注到错误的 token。
实践中，Key 的量化精度通常高于 Value
（Key 需要 INT8，Value 可以 INT4），
因为 Key 直接参与注意力分数的计算，
而 Value 只是被加权求和。

\textbf{动态 KV Cache 驱逐}。
不保存所有历史 token 的 KV cache，
而是维护一个固定大小的缓冲区，
当缓冲区满时驱逐“最不重要”的 KV 对。
重要性的衡量标准通常基于
累积的注意力权重，
如果一个 token 的 KV 长期没有被后续 token 关注，
说明它的信息已经被模型“消化”或不再需要，
可以安全驱逐。
这种策略在某种程度上模拟了人类记忆的“遗忘曲线”，
重要的信息被反复回忆（高注意力权重），
不重要的信息逐渐淡忘（被驱逐）。

\subsection{训练--推理的长度失配}

第三个挑战更加根本：
模型在训练时只见过一定长度的序列，
当推理时序列长度超过训练长度时会发生什么？

RoPE~\cite{su2021rope} 的位置编码依赖于旋转角度
$\theta_i = 10000^{-2i/d}$。
在训练长度内（比如 4K），
模型见过的最大旋转角度是 $4096 \times \theta_i$。
当推理时序列长度变成 32K，
旋转角度变为 $32768 \times \theta_i$，
这些角度值是模型从未在训练中遇到过的。

直觉上，这就像一个只学过 0--100 范围内加法的学生
被要求做 0--1000 范围的加法。
虽然底层规则相同，但具体的数值模式是「未见过」的，
注意力分数变得不可靠，
这就是所谓的「外推失败」（extrapolation failure）。

低频维度受影响最大：
它们的旋转角度在训练范围内变化缓慢（不到一整圈），
但在超长序列中可能转过好几圈，
进入完全陌生的角度空间。
高频维度则相对安全，
因为它们在训练范围内已经转过了很多圈，
多转几圈不会带来本质变化。

我们可以更精确地量化这种“外推危险度”。
对于 RoPE 的第 $i$ 个频率维度，
训练长度 $L$ 内的最大旋转角度是 $L \cdot \theta_i$。
如果 $L \cdot \theta_i \gg 2\pi$（转了很多圈），
那么该维度对任意长度的“新”位置
都有很好的泛化能力，
因为新位置的旋转角度只是已见角度的“自然延伸”。
如果 $L \cdot \theta_i \ll 2\pi$（不到一圈），
那么该维度对超出训练长度的位置几乎没有泛化能力，
新的角度值是模型从未见过的。

以 LLaMA 2 为例（$b = 10000$，$d = 128$，$L = 4096$），
最高频维度（$i = 0$）的训练范围内总旋转角度为
$4096 \times 10000^0 = 4096$ 弧度 $\approx$ 652 圈，
极其安全。
最低频维度（$i = 63$）的总旋转角度为
$4096 \times 10000^{-126/128} \approx 0.04$ 弧度，
连一圈的 $1/150$ 都不到。
这种维度间的巨大差异，
高频维度可以安全地外推到任意长度，
低频维度在训练长度之外完全不可靠，
正是所有位置编码扩展方法需要解决的核心问题。

\subsection{Lost in the Middle：中间信息的遗忘}

第四个挑战来自模型的注意力分配模式。
2023 年 Liu 等人的研究揭示了一个令人不安的现象：
当输入序列很长时，
模型对序列\textbf{开头}和\textbf{结尾}的信息关注度最高，
而对\textbf{中间}部分的信息检索准确度显著下降，
下降幅度超过 30\%。

% NEW REF: liu2023lost = Liu et al., "Lost in the Middle: How Language Models Use Long Contexts", arXiv:2307.03172, 2023

这种「U 型曲线」注意力模式有深层原因。
在自回归训练中，
每个 token 的损失只取决于它之前的 context。
序列开头的 token 是所有后续 token 的 context，
因此它们的信息被反复强化。
序列结尾的 token 距离当前预测位置最近，
在注意力分数中自然占优势（近因效应）。
中间的 token 则两头不沾，
既不是「总是被参考」的锚点，
也不是「最近看到」的新鲜信息。

这个发现意味着：仅仅增加上下文窗口是不够的。
如果关键信息恰好落在长序列的中间部分，
模型可能会「看到但不记住」，
就像一个学生虽然读了整本教科书，
但只记得开头和结尾的内容。

Lost in the Middle 效应的强度
因模型架构和训练方法而异。
一些后续研究发现，
通过在训练数据中刻意将关键信息放在中间位置
（而非总是放在开头或结尾），
可以部分缓解 U 型注意力偏差。
这种“位置增强训练”（positional augmentation）
通过改变训练信号的位置分布，
迫使模型学会对所有位置给予更均匀的关注。

另一种缓解策略是\textbf{双向 attention}（bidirectional attention）
或\textbf{填充式 attention}（fill-in-the-middle）。
在标准的因果注意力中，
每个 token 只能关注它之前的 token，
这种单向性是 U 型偏差的结构性来源之一。
如果在某些层中允许 token 同时关注前后的 token
（类似于 BERT 的双向注意力），
U 型偏差可以被显著降低。
但双向注意力与自回归生成不兼容，
这是一个根本性的架构限制。
iRoPE 的 NoPE 层在某种程度上近似了这种效果：
没有位置编码的注意力层不受因果位置偏置的影响，
使得中间位置的 token 在语义匹配上
不再处于系统性劣势。

\section{位置编码的扩展方法}

面对训练--推理的长度失配问题，
研究者们提出了一系列巧妙的位置编码扩展方法。
这些方法的核心思想都是：
让超长序列中的位置编码落入模型训练时见过的值域范围内。

\subsection{线性位置插值}

最直接的方法是\textbf{线性插值}：
如果要将 4K 上下文扩展到 32K（$8\times$），
就把所有位置索引除以 8。
原来位置 32,768 的 token 变成位置 4,096，
刚好是训练时见过的最大位置。

% NEW REF: chen2023extending = Chen et al., "Extending Context Window of Large Language Models via Positional Interpolation", arXiv:2306.15595, 2023

数学上，将位置 $m$ 替换为 $m/s$，
其中 $s = L_{\text{target}} / L_{\text{train}}$ 是扩展比例。
RoPE 的旋转角度变为 $(m/s) \cdot \theta_i$，
确保所有角度都落在训练时见过的范围内。

为什么这个简单的技巧有效？
因为它保持了\textbf{相对位置关系}不变，
如果 token A 和 token B 的相对位置差是 $\Delta$，
插值后变成 $\Delta / s$，但两者的差值比例没变。
模型学到的「距离越远注意力越弱」的模式仍然适用。

但线性插值有明显的局限：
当扩展比例 $s$ 很大时（比如 $32\times$ 以上），
相邻 token 的位置差变得极小
（从 1.0 变成 $1/32 = 0.03125$），
模型难以区分紧邻的 token 的位置关系，
局部分辨率严重下降。
这就像把一张高清照片压缩到很小：
全局结构保留了，但细节模糊了。

从频域的角度理解线性插值会更加清晰。
RoPE 的每个维度都可以看作一个正弦信号，
频率为 $\theta_i = b^{-2i/d}$。
线性插值将所有频率统一缩放为 $\theta_i / s$，
这等价于将整个频谱向低频方向平移。
问题在于，高频维度（负责编码局部位置关系的维度）
原本就有精细的空间分辨率，
被缩放后分辨率进一步降低，
模型可能连“这个 token 是前一个还是后一个”
这种基本的局部关系都无法区分。
低频维度（编码全局位置的维度）
则被过度压缩，
原本需要几百个 token 才能转完一圈的频率
被压缩到几十个 token 就转完，
全局位置编码的有效区分度下降。

Chen 等人~\cite{chen2023extending}在实验中验证了这一分析：
线性插值在 $4\times$--$8\times$ 的扩展比例下效果良好，
在经过少量微调后（$\sim$1000 步）
困惑度几乎不增加。
但在 $32\times$ 以上的扩展比例下，
即使经过大量微调，困惑度仍然显著增加，
局部分辨率的损失是不可逆的。
这一发现直接启发了后续 NTK-aware 和 YaRN 等方法
对不同频率维度采取差异化处理的思路。

\subsection{NTK-aware 插值}

NTK-aware 插值的核心洞察是：
\textbf{不应该对所有频率维度做相同比例的缩放}。

回忆 RoPE 的频率：$\theta_i = b^{-2i/d}$（$b = 10000$）。
高频维度（$i$ 小）负责编码\textbf{局部}位置关系
（相邻 token 之间的差异），
低频维度（$i$ 大）负责编码\textbf{全局}位置关系
（文档开头 vs 文档结尾）。

线性插值对所有频率一视同仁地除以 $s$，
结果高频维度（负责局部关系的维度）被不必要地压缩了。
NTK-aware 插值只调大 base 值 $b$：
将 $b = 10000$ 替换为 $b' = 10000 \cdot s^{d/(d-2)}$。

这产生了一种非均匀缩放效果：
\begin{itemize}[noitemsep]
  \item \textbf{高频维度}几乎不受影响，
        它们的频率本来就高，稍微调整 base 对它们影响很小。
        局部位置分辨率得以保留。
  \item \textbf{低频维度}被显著缩放，
        让它们能编码更远的位置关系。
\end{itemize}

一个有趣的类比：这就像数字系统中的进制变化。
在十进制中，三位数能表示 0--999。
如果我们把进制从 10 变成 100，
三位数就能表示 0--$100^3 = 999{,}999$，
扩大了 1000 倍的范围，
但低位的分辨率几乎不变。

NTK-aware 插值的最大优势是\textbf{不需要微调}：
在中等扩展比例（$4\times$ 左右）下，
直接修改 base 值就能在推理时扩展上下文，
质量损失很小。但在高扩展比例下仍会退化。

NTK-aware 插值的成功启发了一个更广泛的问题：
\textbf{RoPE 的 base 值到底应该怎么选？}
原始 RoPE 论文~\cite{su2021rope} 使用了 $b = 10000$，
这个值是基于较短序列（$\sim$2K）的实验确定的。
LLaMA 3 在 128K 上下文训练中
使用了 $b = 8 \times 10^6$，
比原始值大了 800 倍。
这个巨大的差异暗示了 base 值的最优选择
强烈依赖于目标上下文长度。

直觉上，更大的 base 值使得所有频率维度的频率更低，
相当于“拉长”了位置编码的周期，
每个维度需要更长的序列才能完成一个完整的旋转周期。
这使得在短序列中相邻位置的区分度降低
（因为旋转角度差变小了），
但在长序列中避免了“角度回绕”
（同一个角度值对应多个不同的位置）。
最优的 base 值需要在短序列分辨率和长序列外推能力之间取得平衡。

一个经验公式是 $b_{\text{opt}} \propto L^{d/(d-2)}$，
其中 $L$ 是目标上下文长度，$d$ 是头维度。
这个公式直接来自 NTK-aware 插值的推导，
它确保了在目标长度 $L$ 下
所有频率维度的外推安全系数至少为 1
（即每个维度在训练范围内至少完成一个完整周期）。

\subsection{YaRN：优雅的频率操控}

YaRN（Yet Another RoPE extensioN）是当前最流行的上下文扩展方法。
它的核心洞察是：RoPE 的不同频率维度应该被\textbf{区别对待}，
而且不仅要调整频率，还要调整注意力的温度。

% NEW REF: peng2023yarn = Peng et al., "YaRN: Efficient Context Window Extension of Large Language Models", arXiv:2309.00071, 2023

RoPE 的频率 $\theta_i = 10000^{-2i/d}$ 从高频（$i$ 小）到低频（$i$ 大）。
YaRN 将这些频率分为三组，并定义了一个关键概念：
频率维度的「波长」$\lambda_i = 2\pi / \theta_i$。
如果 $\lambda_i$ 远大于训练长度 $L$，
说明该维度在训练中连一整圈都没转完，外推风险最高。
如果 $\lambda_i$ 远小于 $L$，
说明该维度已经转过很多圈，外推风险最低。

基于这个分析，YaRN 将频率维度分为三组：
\begin{itemize}[noitemsep]
  \item \textbf{高频维度}（$\lambda_i < L$，对应短距离模式）：
        不做任何修改，这些维度已经能处理局部关系，不需要调整。
        它们在训练中转过了很多圈，
        外推到更长序列时只是多转几圈而已。
  \item \textbf{低频维度}（$\lambda_i > s \cdot L$，对应长距离模式）：
        按上下文扩展比例进行激进的频率缩放，
        比如要将 4K 扩展到 128K（$32\times$），
        这些维度的频率除以 32。
        它们必须被压缩到训练范围内，否则外推完全不可靠。
  \item \textbf{中频维度}（$L \leq \lambda_i \leq s \cdot L$）：
        进行渐进式的 NTK-aware 插值，
        在不缩放和完全缩放之间做平滑过渡。
        过渡函数是一个 ramp（线性插值因子），
        避免了频率空间中的突变。
\end{itemize}

此外，YaRN 引入了\textbf{注意力温度缩放}因子 $\sqrt{t}$。
长序列中注意力分数的方差会增大
（因为更多的 token 参与 softmax 竞争），
导致注意力分布变得过于「尖锐」，
集中在少数 token 上，忽略其余信息。
温度因子 $\sqrt{t}$ 相当于在 softmax 之前
将注意力 logits 除以 $\sqrt{t}$（$t > 1$），
降低 softmax 的「锐度」，
让注意力分布更加平滑，
使模型能更均匀地关注长序列中的信息。

通过这些手段，YaRN 只需在原始预训练数据的极小比例上做微调
（通常 $\sim$0.1\% 的训练数据，训练 $\sim$400 步），
就能将上下文窗口扩展 $32\times$ 甚至更多。
Qwen2.5、DeepSeek 等模型均使用 YaRN 变体进行上下文扩展。

YaRN 的另一个重要贡献是其\textbf{理论分析框架}。
论文将 RoPE 的不同频率维度映射到一个“波长谱”上，
并证明了外推误差与频率维度的波长和训练长度的比值直接相关。
具体来说，对于第 $i$ 个频率维度，
其“外推安全系数” $r_i$ 可以定义为
$r_i = L_{\text{train}} / \lambda_i = L_{\text{train}} \cdot \theta_i / (2\pi)$。
当 $r_i \gg 1$（高频维度，波长远小于训练长度）时，
该维度在训练中已经“见过”了多个完整周期，
外推到更长序列只是多走几个周期，风险极低。
当 $r_i \ll 1$（低频维度，波长远大于训练长度）时，
该维度在训练中连一个周期都没完成，
外推意味着进入完全未见的角度空间，风险极高。

这个安全系数框架为所有基于 RoPE 的位置扩展方法
提供了统一的理论视角。
线性插值对所有维度施加相同的缩放
（无视安全系数的差异），
NTK-aware 通过调整 base 实现了隐式的差异化缩放，
而 YaRN 则显式地按安全系数将维度分为三组，
每组采用不同的处理策略。
后续的各种 RoPE 扩展变体（如 LLaMA 3 使用的频率调整方案）
本质上都是在这个安全系数框架内
寻找不同的分组和缩放策略。

\subsection{Dynamic NTK：推理时自适应}

Dynamic NTK 是一种在推理时根据实际序列长度
\textbf{动态调整}插值比例的方法。

核心思想很简单：不预设固定的扩展比例，
而是在推理时检查当前序列长度 $n$，
如果 $n \leq L_{\text{train}}$，不做任何修改；
如果 $n > L_{\text{train}}$，
将插值比例设为 $s = n / L_{\text{train}}$，
然后应用 NTK-aware 缩放。

这种方法的优势是\textbf{完全免微调}：
不需要任何额外训练，只需修改推理代码。
对于中等程度的扩展（$2\times$--$4\times$），
Dynamic NTK 的质量损失可以接受。
但对于大比例扩展（$8\times$ 以上），
质量下降明显，因为模型没有在长序列上做过任何适应。

Dynamic NTK 在实际应用中常被用作一种「应急方案」，
当用户的输入偶尔超过训练长度时，
它能优雅地降级，而不是直接崩溃。

Dynamic NTK 的一个有趣性质是它的\textbf{非对称效果}。
当序列长度 $n$ 略大于训练长度 $L$ 时
（如 $n = 1.2L$），
Dynamic NTK 几乎不影响前 $L$ 个 token 的表示质量，
因为它们的位置编码值仍在训练范围内，
调整只影响第 $L+1$ 到 $n$ 个 token。
这意味着 Dynamic NTK 可以被看作一种
“保守的超限处理”，
已有 token 的质量不受影响，
超限 token 获得“尽力而为”的处理。

多个开源推理框架（如 vLLM、SGLang、llama.cpp）
都内置了 Dynamic NTK 作为默认的超限处理策略。
当推理请求的长度超过模型的训练上下文窗口时，
框架自动启用 Dynamic NTK 而不是拒绝请求或截断输入。
这种“无感降级”的用户体验设计
使得 Dynamic NTK 成为了实际部署中使用最广泛的
位置编码扩展方法，
尽管它在极端扩展比例下效果不如 YaRN。

\subsection{训练侧正则化：Shuffle the Context 与 Phase-wise RoPE 缩放}
\label{sec:rope-train-reg}

PI、YaRN、LongRoPE 这条线本质上是\textbf{推理时的频率操控}，
不改变训练目标。
2026 年第二季度出现的两项工作
把视角从推理侧转回\textbf{训练侧}，
通过正则化和阶段化训练
让模型本身就具备更好的长度外推能力。

\textbf{Shuffle the Context}（RoPE-Perturbed Self-Distillation）
~\cite{shuffle2026}（arXiv:2604.14339，Microsoft，2026-04-15）
的关键 trick 是\textbf{用 RoPE 扰动构造两个视图}：
对同一段长文本，分别用原始 RoPE 和扰动后的 RoPE
（如随机偏移 base 频率、或对部分维度做 shuffle）
跑两次前向，
然后用自蒸馏 loss 强迫两个视图的输出 logits 一致。
这个设计的直觉是：
\textbf{如果模型真正"理解"了内容，
那么任何合理的位置编码扰动都不应改变它的回答}。
扰动一致性约束让模型把推理能力从特定 RoPE 角度中解耦。
实测：Llama-3-8B 在 RULER-64K 上
\textbf{+12.04\%}；Qwen-3-4B 在 RULER-256K 上 \textbf{+2.71\%}。

\textbf{Short Data, Long Context}~\cite{shortdatalongctx2026}
（arXiv:2604.06070，Meta，2026-04-07）
解决的是另一个工程痛点：
\textbf{真长样本稀缺且昂贵}。
长文档训练数据的获取成本远高于短样本，
且大部分长样本是低质量拼接。
论文提出 \textbf{Phase-wise RoPE scaling}——
每个训练阶段最大化对应频率段的 rotational spectrum 利用率，
配合 logit-based KD（用强短上下文模型作 teacher）
\textbf{仅用打包短样本}就能迁移长上下文检索能力。
论文还首次给出可视化：
位置扰动如何通过 q/k → 后续 transformer layer
→ output logits 系统传播。
这一发现让"短数据训长上下文"
从经验技巧变为可量化的训练策略。

两条工作给 RoPE 外推带来一个新的思路框架：
PI/YaRN/LongRoPE 通过\textbf{推理时频率重映射}让模型勉强外推，
Shuffle/Short-Data 通过\textbf{训练时正则化}
让模型从一开始就对长度具有鲁棒性。
两者正交可叠加。

\section{iRoPE：Meta 的交替架构}

LLaMA 4 引入了更激进的方案，iRoPE（interleaved RoPE）。
它不是在位置编码的数值上做文章，
而是从\textbf{架构层面}重新思考位置编码的角色。

\subsection{核心设计：交替堆叠}

iRoPE 交替堆叠两种注意力层：
\begin{itemize}[noitemsep]
  \item \textbf{RoPE 层}：使用旋转位置编码，
        捕捉局部上下文和精确的相对位置关系。
        这些层知道「哪个 token 在前，哪个在后」。
  \item \textbf{NoPE 层}（No Positional Encoding）：
        完全没有位置编码，注意力只基于语义内容。
        这些层天然不受序列长度限制，
        因为它们根本不编码位置信息。
\end{itemize}

为什么这种交替架构有效？
关键洞察是：\textbf{不是所有注意力模式都需要位置信息}。

考虑两种典型的注意力模式：
一种是「这个代词指代的是三句话前的名词」：
这需要精确的位置感知（RoPE 层的工作）。
另一种是「这段话在讨论量子物理的概念」，
这只需要语义匹配，不关心具体位置（NoPE 层的工作）。

NoPE 层的存在使得模型对长序列有天然的适应性，
它们关注的是「这个 token 的内容是什么」
而非「这个 token 在哪个位置」。
当序列从 8K 扩展到 1M 时，
NoPE 层完全不受影响，
只有 RoPE 层需要做位置编码的扩展。
这等于把需要外推的层数减少了一半。

iRoPE 的设计理念有一个深刻的理论基础。
研究者观察到 Transformer 中的注意力模式
可以大致分为两类：
\textbf{位置敏感型}（positional）和\textbf{内容敏感型}（content-based）。
位置敏感型注意力关注的是“这个 token 距离我有多远”，
典型的例子是局部语法模式
（如“形容词通常在名词前一两个位置”）。
内容敏感型注意力关注的是“这个 token 和我的语义有多相关”，
典型的例子是指代消解
（代词需要找到它指代的名词，
不管该名词在多远的位置）。

在纯 RoPE 架构中，
两类注意力模式都被“强加”了位置编码，
即使是纯内容匹配的注意力头，
也会被 RoPE 的旋转角度所影响。
这在短序列中不是问题
（因为旋转角度变化不大），
但在超长序列中，RoPE 的角度值偏离了训练分布，
内容匹配的质量被位置编码的噪声所干扰。

iRoPE 通过为两类注意力分配专门的层
（RoPE 层处理位置敏感型，NoPE 层处理内容敏感型），
让每种注意力模式在最适合的计算环境中运行。
这种“关注点分离”（separation of concerns）
不仅改善了长上下文的外推能力，
还在短序列上带来了轻微的质量提升，
因为内容匹配不再受不必要的位置偏置干扰。

\subsection{LLaMA 4 的 iRoPE 实现细节}

LLaMA 4 系列包含三个模型，
全部采用 iRoPE + MoE 架构：

\begin{center}\small
\renewcommand{\arraystretch}{1.3}
\begin{tabular}{lccccc}
\toprule
\rowcolor{figLightGray}
\textbf{模型} & \textbf{总参数} & \textbf{活跃参数} & \textbf{专家数} & \textbf{上下文} \\
\midrule
Scout   & 109B  & 17B  & 16    & 10M tokens \\
Maverick & 400B & 17B  & 128   & 1M tokens \\
Behemoth & 2T   & 288B & 16    & --- (训练中) \\
\bottomrule
\end{tabular}
\end{center}

Scout 的 10M token 上下文窗口（约 15,000 页文本）
是截至 2025 年 4 月所有公开模型中最长的。
这个惊人的长度正是得益于 iRoPE 的架构设计：
NoPE 层提供了「免费」的上下文扩展空间，
RoPE 层通过频率调整处理位置敏感的注意力模式。

在 LLaMA 4 的具体实现中~\cite{meta2025llama4,llama4irope2026}，
iRoPE 的交替模式并非简单的 1:1 交替。
Meta 发现不同的交替比例适合不同的模型规模，
Scout 和 Maverick 使用了针对各自专家数量和模型规模
优化的交替策略。
Scout 的具体配置是\textbf{每 4 层中前 3 层为 RoPE + chunked attention}
（每个 token 仅看本地 8192 token 的窗口），
\textbf{第 4 层为 NoPE + full causal attention}（看完整序列），
循环堆叠——75\% 的层只看 8K 局部窗口，
25\% 的层周期性扫全文。
为防止极长序列下 softmax 分布塌缩，
NoPE 层在推理时采用 \textbf{Scalable Softmax}
（SSMax）~\cite{nakanishi2025ssmax} 的温度缩放策略，
让注意力权重在 1M+ 序列长度下仍保持合理的尖锐度。
这一组合（NoPE + SSMax + chunked RoPE）
将 NoPE 论文（Kazemnejad 2023）~\cite{kazemnejad2023nope}
从研究 curiosity 推到了工业模型的首个大规模实现。

iRoPE 还与 LLaMA 4 的其他架构创新协同工作：
\begin{itemize}[noitemsep]
  \item \textbf{早期融合}（Early Fusion）：
        LLaMA 4 是原生多模态模型，
        文本、图像和视频 token 在模型早期就被融合处理。
        NoPE 层对跨模态注意力特别有价值，
        「这张图片展示的概念」和「文本描述的概念」
        之间的匹配是纯语义的，不需要位置信息
  \item \textbf{MoE 稀疏激活}：
        每个 token 只激活少量专家，
        使得 Scout 的 109B 总参数中仅有 17B 活跃，
        这让 10M 上下文在 8$\times$H100 上可运行
\end{itemize}

\begin{keybox}[iRoPE 的训练流程]
LLaMA 4 的长上下文训练分三个阶段：
\begin{enumerate}[noitemsep]
  \item \textbf{预训练}：在 8K 上下文长度上完成主体训练。
        这个阶段消耗了绝大部分计算。
        Behemoth 作为「教师模型」在这个阶段训练，
        Scout 和 Maverick 通过蒸馏获得高质量初始化。
  \item \textbf{中等长度扩展}：调整 RoPE 层的频率参数，
        将上下文扩展到 128K--256K，
        在长文档数据上做少量继续训练。
  \item \textbf{超长扩展}：进一步扩展到 1M--10M。
        NoPE 层在这个阶段几乎不需要适应，
        RoPE 层通过 YaRN 风格的频率操控进行扩展。
\end{enumerate}
LLaMA 4 Scout 通过 iRoPE 达到了 10M token 的上下文窗口，
这在纯 RoPE 架构中几乎不可能实现。
\end{keybox}

\subsection{iRoPE 的推理优化}

iRoPE 的另一个优势是推理效率。
NoPE 层可以使用更激进的 KV cache 压缩策略
（比如丢弃远距离的 KV），
因为它们的注意力模式基于语义相似度而非位置，
远距离的 token 只要语义不相关就不重要。
而 RoPE 层则保留完整的位置信息，
确保精确的位置感知不受影响。

这种差异化的 KV cache 管理策略
在实践中带来了显著的内存节省。
对于 NoPE 层，可以采用\textbf{语义驱逐}策略，
当 KV cache 接近容量上限时，
驱逐与当前 query 语义相似度最低的 KV 对。
由于 NoPE 层不编码位置信息，
被驱逐的 KV 对只要语义不相关就不会影响注意力质量。
而对于 RoPE 层，则需要采用更保守的
\textbf{窗口驱逐}策略（保留最近的 $W$ 个 token 和开头的 $H$ 个 token），
因为位置信息的丢失可能导致位置感知的退化。
综合两种策略，iRoPE 模型在 10M 上下文下
KV cache 的实际内存占用
可以降低到全量存储的约 40\%--60\%，
这是纯 RoPE 架构难以达到的压缩比。

vLLM 在 v0.8.3 版本中已原生支持 LLaMA 4 的 iRoPE 架构。
实测显示，使用 8$\times$H100，
vLLM 可以为 Scout 提供 1M 上下文的推理服务，
为 Maverick 提供约 430K 上下文的推理服务。
值得注意的是，Scout 的标称上下文为 10M，
但当前 vLLM 实现只能服务到 1M，
差距来自 KV cache 的显存限制。
要在 8$\times$H100 上支持完整的 10M 上下文，
需要结合上述的差异化 KV cache 压缩
和系统级的内存管理优化，
这仍是活跃的工程研究方向。

\begin{warnbox}[iRoPE 的局限与争议]
LLaMA 4 在发布后收到了一些社区反馈：
虽然 10M token 的上下文窗口在 NIAH 等检索测试上表现优秀，
但在需要\textbf{深度推理}的长上下文任务上，
实际有效利用的上下文长度远低于标称值。
2026 年的多项独立评测显示~\cite{llama4irope2026}，
Scout 在大约 \textbf{300K token} 的综合任务上即出现明显退化，
所谓“10M 上下文”更多是 NIAH 级别的“硬塞得下”，
远未达到“读得懂”。
这再次验证了「标称上下文长度 $\neq$ 有效上下文长度」的原则。
此外，Behemoth 模型截至 2026 年初仍未正式发布，
社区对其实际性能持审慎态度。
\end{warnbox}

\section{NoPE 路线：从研究 curiosity 到主流配置}
\label{sec:nope-route}

NoPE 在 2023 年提出时~\cite{kazemnejad2023nope}
被视为一个反直觉的研究 curiosity：
Kazemnejad 等证明了
\textbf{decoder-only 模型仅靠 causal mask 就能学到位置},
不需要任何显式的位置编码（不需要 RoPE，不需要 ALiBi，不需要任何东西）。
原因在于 causal mask 本身就携带顺序信息——
位置 $t$ 的 token 知道自己只能看见 $1$ 到 $t-1$ 的 token，
这种"可见性梯度"足以让自注意力学到隐式的相对位置。

2025--2026 年的工业实践证明，
NoPE 不仅可行，而且在长上下文外推上反而\textbf{更优}。
原因是显式位置编码（如 RoPE）在外推到训练分布之外时
会产生角度漂移，
而 NoPE 没有这种漂移源。
两个旗舰应用同时把 NoPE 推向主流：
\begin{itemize}[noitemsep]
  \item \textbf{Llama 4 Scout/Maverick 的 iRoPE}~\cite{llama4irope2026}：
        每 4 层中第 4 层用 NoPE + full causal attention，
        前 3 层用 RoPE + chunked attention，
        在 8K 局部 + 周期性全局之间寻得平衡
  \item \textbf{Mamba-3 hybrid}~\cite{mamba3hybrid2026}：
        每 5 层 SSM 插 1 层 NoPE Self-Attention。
        SSM 的递归状态已经隐式编码顺序，
        attention 层只需做内容匹配，不需要位置信号
\end{itemize}

NoPE 在长序列上的一个工程问题是 softmax 塌缩：
随序列长度增加，注意力权重被稀释到接近均匀，
任何一个 token 都拿不到足够的注意力质量。
\textbf{Scalable Softmax}（SSMax，Nakanishi 2025）
~\cite{nakanishi2025ssmax}
通过推理时温度缩放
（softmax 分母乘以 $\log_n s$ 形式的因子）
把注意力分布重新拉回尖锐区间，
成为 NoPE 在 1M+ 序列上的标准稳定剂。
SSMax 的存在让"NoPE 在 10M 序列上仍能正常 attend"
从理论可能变为生产现实。

NoPE 的产业意义远不止"少做一件事"。
它把\textbf{位置编码的扩展难题}（PI、YaRN、LongRoPE 这条线）
从工业模型的关键路径上移除：
只要 NoPE 层占足够比例，
模型就天然支持任意长度外推，
不再需要为每个新长度调整 RoPE base 频率。
2026 年初，多个前沿团队的内部实验已证实
\textbf{纯 NoPE 也可训}（小规模上 NoPE 与 RoPE 性能持平甚至更优），
唯一限制是缺少大规模 SOTA 验证——
而 Llama 4 + Mamba-3 已经从两个不同方向给出了肯定答案。

\section{DeepSeek 稀疏注意力演进：NSA $\to$ MoBA $\to$ DSA $\to$ NSMLA $\to$ V4}
\label{sec:deepseek-sparse-attention}

如果说 RoPE 外推是长上下文的"位置编码主线"，
那么稀疏注意力就是\textbf{算力主线}：
注意力的 $O(n^2)$ 是 KV cache 之外的另一只"长上下文野兽"，
1M token 在密集 attention 上意味着 $10^{12}$ 次 dot product。
2025 年初到 2026 年中，DeepSeek 沿着稀疏注意力一条线走了 5 步，
从 NSA 的"训练-推理一致稀疏"
到 V4 的"CSA + HCA + SWA 三层混合"，
把可训练稀疏注意力从论文推进到旗舰生产。

\textbf{NSA}（Native Sparse Attention，arXiv:2502.11089）
~\cite{nsa2025} 是 2025 年 2 月的奠基之作，
也是首个在大规模训练中证明
\textbf{稀疏注意力可以端到端训练}的工作。
NSA 的设计三件套：
（1）\textbf{coarse-grained token 压缩}——
用一个轻量小网络把每个块的 token 压缩成单一表示，
全局 attention 在压缩表示上做；
（2）\textbf{块级选择}——
对每个 query，从压缩表示中选 top-$k$ 块，
把这些块还原为原始 token 参与精细 attention；
（3）\textbf{sliding window}——
保证局部强连接不丢失。
NSA 在 27B/260B token 训练规模下，
64K 序列 forward 加速 \textbf{9$\times$}、backward 加速 \textbf{6$\times$}，
质量持平甚至超过 full attention，
是 ACL 2025 最佳长论文之一。

\textbf{MoBA}（Mixture of Block Attention，arXiv:2502.13189）
~\cite{moba2025} 几乎同期由 Moonshot 提出。
MoBA 的洞察是把\textbf{MoE 的 gating 思想搬进 attention}：
对每个 query，用 parameter-less 的 top-$k$ gating
从 KV 块中选最相关的几块。
关键是 gating 完全不引入新参数，
只是对块内 key 的均值与 query 做点积排序，
"接管"了 NSA 的块级选择步骤但更轻量。
Moonshot 已把 MoBA 用在 Kimi 的生产长上下文请求上；
基于 Llama-3.1-8B 的 MoBA 长上下文版本
在 100B token 激活规模下达到 \textbf{95.31\%} 的注意力稀疏度，
仅用 5\% 的 attention 计算就能保持质量。

\textbf{DSA}（DeepSeek Sparse Attention，V3.2）
~\cite{deepseekv32dsa} 在 2025 年 9 月把可训练稀疏 attention
推进到\textbf{产品级}。
DSA 的核心是 \textbf{Lightning Indexer}：
一个 FP8 精度的少头小网络，
对历史 token 打分，
然后在每个 query 上做 \textbf{Fine-Grained Top-2048 KV} 选择。
这把每个 query 的 KV 访问从 $O(L)$ 降到 $O(k=2048)$，
注意力总复杂度从 $O(L^2)$ 降为 $O(Lk)$。
vLLM 在 V3.2 发布日（2025-09-29）即原生支持：
\textbf{prefill 与 decode 分别处理 indexer}，
paged attention 缓存布局相应分离；
DeepGEMM 提供 \texttt{fp8\_mqa} 和 paged indexer kernel；
FlashMLA 提供稀疏 attention kernel。
工程坑在于 indexer 与 MLA 用不同的 RoPE layout，
任何同时使用 DSA 和 MLA 的实现都必须同步两套位置编码。

\textbf{NSMLA}（OpenReview r9xDZJgIKk，2026 ICLR/NeurIPS 评审中）
~\cite{nsmla2026} 是 DSA 的训练范式补全。
DSA 的局限是 lightning indexer 通过\textbf{蒸馏}从 full attention 训练出来——
意味着仍然需要完整 attention 的预训练阶段。
NSMLA 改造 DSA 让其\textbf{从零可预训练}：
indexer 在原生形态下与 MLA 联合训练，
推理时再转换为 lightning 形态。
工程实现使用 TileLang 写训练 kernel、
复用 vLLM 推理基础设施。
正是 NSMLA 把训练成本一次性降下来后，
DeepSeek 在 2026 年 1 月把 \textbf{API 价格降低 ≥50\%}，
首次让百万 token 服务化在经济上变得可行。

\textbf{DeepSeek-V4}~\cite{deepseekv4}（2026-04-24）
是这条主线的当前终点：
V4-Pro（1.6T 总参 / 49B 活跃）
和 V4-Flash（284B 总参 / 13B 活跃）
原生支持 1M 上下文，
注意力被设计为\textbf{三种层间交错}：
\begin{itemize}[noitemsep]
  \item \textbf{CSA}（Compressed Selective Attention）：
        先沿序列维压缩 KV，
        再在压缩表示上做 top-$k$ 稀疏选择
  \item \textbf{HCA}（Hierarchical Compressed Attention）：
        更激进的全局压缩，
        牺牲局部精度换全局视野
  \item \textbf{SWA}（Sliding Window Attention）：
        保证最近 token 的强连接，
        弥补 CSA/HCA 的局部稀疏
\end{itemize}
V4 vs V3.2 在 1M 上下文：
\textbf{per-token 推理 FLOPs 降低 73\%、KV cache 降低 90\%}。
配合 Muon 优化器和 on-policy 蒸馏后训练，
V4 是首个把"1M 上下文 = 经济默认"
落到\textbf{开源旗舰}的模型。
DeepSeek 自评 V4-Pro-Max（扩展 reasoning tokens 版本）
"trails state-of-the-art frontier models by approximately 3 to 6 months"——
首次官方量化开源-闭源差距。

\begin{keybox}[NSA $\to$ MoBA $\to$ DSA $\to$ NSMLA $\to$ V4 的工程线索]
\begin{itemize}[noitemsep]
  \item \textbf{NSA}：证明 trainable sparse attention 在大规模上可行
  \item \textbf{MoBA}：用 MoE-style gating 替代 NSA 的块选择，更轻量
  \item \textbf{DSA}：把 NSA 思路 + lightning indexer 推到产品级，vLLM Day-0 支持
  \item \textbf{NSMLA}：让 indexer 从零可预训练，API 价格再降 ≥50\%
  \item \textbf{V4 (CSA+HCA+SWA)}：三层混合，1M 上下文 -73\% FLOPs / -90\% KV
\end{itemize}
两条主线对照：
NSA / DSA / V4 走"分层压缩 + 选择"路线，
MoBA 走"MoE-style gating"路线。
两者在 V4 中实质性融合（CSA = 压缩 + 选择 + gating）。
\end{keybox}

\section{上下文压缩：用更少的 token 表达更多的信息}
\label{sec:context-compression}

一种完全不同的长上下文策略是：
\textbf{不扩展上下文窗口，而是压缩放入上下文的内容}。
如果我们能在不丢失关键信息的前提下
将 100K token 的文档压缩成 20K token，
那么一个 32K 上下文的模型就足以处理它，
不需要任何位置编码扩展或稀疏注意力优化。

\subsection{LLMLingua：基于模型的 Prompt 压缩}

LLMLingua（2024）及其后续版本 LongLLMLingua% NEW REF: llmlingua2024 = LLMLingua: Compressing Prompts for Accelerated Inference, Jiang et al., EMNLP 2024
提出了一种\textbf{基于困惑度}的 prompt 压缩方法。
核心思想是使用一个小型语言模型
评估 prompt 中每个 token 的“信息量”，
困惑度高的 token（对模型来说意外的、信息量大的 token）
被保留，困惑度低的 token（可预测的、冗余的 token）
被删除或替换为更简洁的表达。

LLMLingua 的压缩流程分为三个步骤：
\begin{enumerate}[noitemsep]
  \item \textbf{粗粒度筛选}：将 prompt 分成段落级的块，
        计算每个块对最终回答的预期贡献度，
        删除最不相关的块。
  \item \textbf{细粒度压缩}：在保留的块内部，
        逐 token 评估困惑度，
        删除可预测性最高（困惑度最低）的 token。
  \item \textbf{语法修复}：对删除 token 后可能产生的
        语法不完整进行后处理修复。
\end{enumerate}

LLMLingua 可以将 prompt 压缩到原始长度的 10\%--30\%，
同时在大多数下游任务上保持 95\%+ 的原始性能。
这意味着一个 128K 上下文的模型
通过 LLMLingua 可以“等效”处理 400K--1.2M token 的输入，
性价比远高于直接扩展上下文窗口。

但 LLMLingua 有一个根本局限：
\textbf{它是一种有损压缩}。
被删除的 token 可能包含在特定查询下变得关键的信息，
一个在当前看起来不重要的细节，
可能恰好是用户即将提问的目标。
这种“无法预知未来查询”的困境
使得 LLMLingua 更适合
已知查询意图的场景（如 RAG 中的检索结果压缩），
而不太适合开放式的长文档探索。

\subsection{StreamingLLM：无限流式推理}

StreamingLLM~\cite{xiao2024streamingllm}（2024）% NEW REF: xiao2024streamingllm = Xiao et al., "Efficient Streaming Language Models with Attention Sinks", arXiv:2309.17453, 2024
发现了一个关于注意力机制的重要现象：
\textbf{注意力沉降}（attention sinks）。
在自回归语言模型中，
序列的前几个 token（通常是前 1--4 个）
会“吸收”大量的注意力权重，
即使这些 token 的语义内容并不重要。

这个发现的直觉解释是：
softmax 注意力强制所有注意力权重之和为 1。
当一个 query 与所有 key 都不太相关时，
注意力权重需要分配到某些地方，
前几个 token 因为在所有后续 token 的注意力计算中
都作为候选存在，逐渐被“训练”成
默认的注意力“倾倒点”。
它们的角色不是提供有用的语义信息，
而是充当 softmax 归一化的“缓冲区”。

StreamingLLM 利用这一发现实现了\textbf{无限长度}的流式推理：
只保留前几个“attention sink” token 的 KV cache
和最近 $W$ 个 token 的 KV cache（$W$ 是窗口大小），
丢弃中间的所有 KV cache。
这将 KV cache 的大小固定为 $O(W)$
而非 $O(n)$（$n$ 为序列长度），
使得模型可以在固定内存下处理无限长的输入流。

StreamingLLM 的代价是显而易见的：
窗口之外的信息被\textbf{永久丢弃}。
它不是一种真正的长上下文方案，
而是一种“无限短上下文”方案：
模型只能“记住”最近的 $W$ 个 token
和最初的几个 attention sink token。
但在某些特定场景下（如实时对话、流式日志分析），
这种“只看最近”的策略已经足够好。

\subsection{AutoCompressor：学习压缩记忆}

AutoCompressor（2023）% NEW REF: autocompressor2023 = AutoCompressor: Learning to Compress Long Contexts, Chevalier et al., arXiv:2305.14788, 2023
代表了另一种压缩哲学：
训练模型\textbf{自己学会}如何将长上下文压缩成少量的“摘要 token”。
模型在处理完一段文本后，
将其压缩为固定数量（如 50 个）的“软摘要”向量，
这些向量作为后续段落的额外上下文输入。

这种方法的优雅之处在于：
压缩是端到端学习的：
模型自主决定哪些信息值得保留，
哪些可以丢弃。
通过在训练中最大化下游任务（如问答、摘要）的性能，
模型学会了一种任务最优的压缩策略。

但 AutoCompressor 面临一个“信息瓶颈”问题：
50 个摘要 token 能编码多少信息？
如果原始文本有 100K token，
压缩比为 2000:1，
即使每个摘要 token 是 4096 维的向量，
总信息容量也远低于原始文本。
实验表明，AutoCompressor 在保持对话一致性和
主题跟踪等“粗粒度”任务上表现良好，
但在需要精确引用原文细节的任务上
性能显著下降。

\subsection{Landmark Attention：标记重要位置}

Landmark Attention~\cite{mohtashami2023landmark}（2023）% NEW REF: mohtashami2023landmark = Mohtashami and Jaggi, "Landmark Attention: Random-Access Infinite Context Length for Transformers", arXiv:2305.16300, 2023
引入了一种“地标”机制：
在文本中插入特殊的 landmark token，
这些 token 作为周围文本块的“代表”。
在注意力计算中，
远距离的 token 只与 landmark token 交互
（而非逐个 token 交互），
只有当 landmark token 的注意力分数足够高时，
才“展开”对应文本块中的所有 token 参与精细计算。

这种设计本质上是一种\textbf{分层注意力}，
粗粒度层级使用 landmark 快速筛选相关区域，
精细粒度层级只在相关区域内做全注意力。
与 Double-P~\cite{doublep2026} 的分层 top-$p$ 策略有异曲同工之处，
但 Landmark Attention 的“地标”是在训练中学习的，
而 Double-P 的分块筛选是基于注意力分数动态计算的。

Landmark Attention 的一个独特优势是它支持\textbf{随机访问}，
模型可以直接“跳转”到任意 landmark 位置，
不需要顺序扫描整个序列。
这使得它特别适合
需要在长文档中精确定位特定信息的任务
（如法律文档中的条款检索），
而传统的注意力机制必须遍历所有位置才能完成这种检索。

\subsection{2026 年的 KV-Cache 压缩：从散点压缩到时间分层}
\label{sec:kv-compression-2026}

2026 年第一季度，KV-cache 压缩一次性出了四篇高质量工作，
覆盖了从\textbf{散点压缩}（DASH-KV、TriAttention）
到\textbf{结构感知}（LCA）
再到\textbf{时间分层}（TTKV）的完整谱系。
这些方法共同把"长上下文 = KV 爆炸"
从架构问题变为可分层调优的工程问题。

\textbf{DASH-KV}~\cite{dashkv2026}（arXiv:2604.19351，2026-04-21）
把注意力\textbf{重新表述为 KV 的 asymmetric deep hashing
近邻搜索}：
Q 和 K 各自学一个非对称的哈希函数，
在哈希空间中对每个 query 只检索 top-$k$ 邻居。
关键 token 仍以全精度回退保护，
其余在低精度哈希空间内交互。
LongBench 上 DASH-KV 匹配 full attention，
但把 attention 复杂度从 $O(N^2)$ 降到 \textbf{$O(N)$}——
首次让线性 attention 在生产场景下与 full attention 等质量。

\textbf{TriAttention}~\cite{triattention2026}（arXiv:2604.04921，2026-04-06）
源于一个被忽略的观察：
\textbf{pre-RoPE 的 Q/K 高度集中在固定非零中心}。
这个观察意味着 key 之间的"语义距离"
可以用三角级数（trigonometric series）排序。
TriAttention 用三角级数为每个 key 算重要性分数，
在 AIME25 32K 生成下达到
\textbf{与 Full Attention 持平的精度，
吞吐 2.5$\times$ 或 KV 显存减少 10.7$\times$}（二选一）。

\textbf{LCA}（Latent Context Anchoring）~\cite{lcakv2026}
（arXiv:2604.12452）
专门针对 \textbf{MLA latent space} 设计：
在 MLA 的低秩潜空间内做 context condensation，
辅以 anchor 位置 key 保留绝对定位能力。
LCA 在 128K 上下文上实现
\textbf{prefill 加速 2.5$\times$、KV cache 减少 90\%}，
是与 DeepSeek MLA 架构最契合的压缩方案。

\textbf{TTKV}~\cite{ttkv2026}（arXiv:2604.19769，2026-03-27）
引入了 KV cache 的\textbf{时间分层}思想：
近期 token 的 KV 留在 HBM 高精度，
远期 token 的 KV 下沉到 DRAM 低精度，
访问时按需提升精度。
128K 任务上的实测：
\textbf{跨层流量降低 5.94$\times$、延迟降低 76\%、吞吐 2$\times$}。
TTKV 把 KV cache 管理从纯空间问题
扩展到\textbf{时间-空间联合调度}，
是 2026 年 KV cache 工程的新轴。

\begin{keybox}[KV-Cache 压缩方法对照（2026 年）]
\begin{itemize}[noitemsep]
  \item \textbf{DASH-KV}：散点压缩 + 哈希近邻，$O(N)$ 等质量
  \item \textbf{TriAttention}：pre-RoPE 集中性 + 三角级数排序，2.5$\times$ 吞吐
  \item \textbf{LCA}：MLA 潜空间专用，128K prefill 2.5$\times$、KV -90\%
  \item \textbf{TTKV}：时间分层，HBM/DRAM 协同，延迟 -76\%
\end{itemize}
四者正交，可叠加使用。
\end{keybox}

\section{Ring Attention：跨 GPU 处理超长序列}

即使解决了位置编码的问题，
超长序列的 KV cache 仍然可能超过单块 GPU 的内存容量。
Ring Attention 的解决方案是
将长序列\textbf{分割到多块 GPU 上}。

\subsection{核心思想：环形传递}

Ring Attention 的灵感来自一个简单的观察：
标准注意力要求每个 query 都能访问所有的 key 和 value，
但没人说所有 KV 必须\textbf{同时}在同一块 GPU 上。

设备排列成环形拓扑。
每块 GPU 持有序列的一个 chunk（包含该 chunk 的 Q、K、V），
然后执行以下步骤：
\begin{enumerate}[noitemsep]
  \item 每块 GPU 计算自己持有的 Q 与\textbf{当前本地} KV 的注意力。
  \item \textbf{同时}，每块 GPU 将自己的 KV block
        发送给环上的下一块 GPU，
        并从前一块 GPU 接收新的 KV block。
  \item 收到新的 KV block 后，计算 Q 与这些新 KV 的注意力，
        将结果与之前的注意力输出做\textbf{在线聚合}
        （online softmax aggregation，
        不需要存储完整的注意力矩阵，
        只需要维护 running sum 和 running max）。
  \item 重复步骤 2--3，直到 KV block 在环上转了一整圈，
        每块 GPU 都看到了所有的 KV。
\end{enumerate}

关键的工程技巧在于步骤 2：
\textbf{通信与计算重叠}。
当 GPU 在计算当前 KV block 的注意力时，
下一个 KV block 正在通过网络传输过来。
只要计算时间大于传输时间，
在序列足够长的情况下通常成立，
传输延迟就被完全隐藏了。

\subsection{内存与扩展性分析}

Ring Attention 将每块 GPU 的内存需求
从 $O(n)$（存储完整 KV cache）降为 $O(n/P)$
（$P$ 是 GPU 数量），
理论上支持无限长的序列，
只要有足够多的 GPU。
PyTorch 已经将 Ring Attention 原生集成到
ContextParallel API 中。

但实际的扩展性受两个因素限制：
\begin{itemize}[noitemsep]
  \item \textbf{通信带宽}：每个 KV block 需要在环上传递 $P-1$ 次。
        当 GPU 数量很多时，KV 的总通信量为 $O(n \cdot (P-1))$，
        网络带宽成为瓶颈。
        在 NVLink（900 GB/s）互连的节点内这不是问题，
        但跨节点（通常 100--400 Gb/s InfiniBand）时延迟显著增加。
  \item \textbf{收益递减}：超过一定长度后，
        模型的有效信息利用率下降（context rot 效应），
        更多的 GPU 和更长的上下文不一定带来更好的结果。
\end{itemize}

\subsection{Zig-Zag Ring Attention：负载均衡}

Ring Attention 有一个变体叫 \textbf{Zig-Zag Ring Attention}，
它解决了原始方法在自回归模型中的负载不均衡问题。
在因果注意力中（每个 token 只能关注之前的 token），
序列开头的 GPU 计算量最小（只有少量的 KV 对），
序列末尾的 GPU 计算量最大（需要关注所有之前的 token）。

如果将序列按顺序均匀分配给 $P$ 块 GPU，
第 1 块 GPU 的计算量约为第 $P$ 块的 $1/P$，
当 $P = 8$ 时，最轻的 GPU 只有最重的 $1/8$ 的工作量，
7/8 的时间在空等。

Zig-Zag 通过锯齿形的序列分割策略平衡各 GPU 的负载：
第一块 GPU 拿第 1、4、5、8 段序列，
第二块拿第 2、3、6、7 段……
使得每块 GPU 处理的总 token 数近似相等。
直觉上，每块 GPU 同时拥有序列开头（计算量小）
和序列末尾（计算量大）的 chunk，
两者的计算量互相平衡。

\textbf{Striped Attention} 是另一种变体，
它将序列按 token 级别交错分配
（第 $i$ 个 token 分配给第 $i \mod P$ 块 GPU），
实现了理论上完美的负载均衡，
但通信模式更复杂。

\subsection{Ring Attention 与 FlashAttention 的结合}

Ring Attention 和 FlashAttention 分别解决了
长序列注意力的两个不同瓶颈：
Ring Attention 将 KV cache \textbf{分布在多块 GPU 之间}
以解决单 GPU 内存不足的问题，
FlashAttention 通过 \textbf{分块计算和在线 softmax}
减少单块 GPU 上的 HBM 读写量。

两者的结合是自然的：
在 Ring Attention 的每一步中，
当 GPU 收到一个新的 KV block 后，
使用 FlashAttention 的分块算法
计算本地 Q 与该 KV block 之间的注意力，
并使用在线 softmax 进行增量聚合。
FlashAttention 的在线 softmax
正好为 Ring Attention 的增量注意力计算
提供了数学基础，
它允许在不保存完整注意力矩阵的前提下
逐块累积注意力输出。

PyTorch 的 ContextParallel API 在底层
正是采用了这种 Ring Attention + FlashAttention 的组合。
用户只需指定上下文并行度 $P$
（将序列拆分到 $P$ 块 GPU 上），
框架自动处理环形通信、分块注意力计算
和在线 softmax 聚合。
这种“一键式”的长上下文训练支持
极大地降低了超长序列训练的工程门槛。

\subsection{实际吞吐量与扩展效率}

Ring Attention 的实际扩展效率
受到计算--通信比（compute-communication ratio）的制约。
假设每块 GPU 上有 $n/P$ 个 token（$n$ 为总序列长度，$P$ 为 GPU 数量），
每一步的计算量约为 $O((n/P)^2 \cdot d)$（本地 Q 与 KV block 的注意力），
通信量为 $O((n/P) \cdot d)$（传输一个 KV block）。
计算--通信比为 $O(n/P)$，
当 $n/P$ 足够大时（每块 GPU 上有足够多的 token），
计算时间远大于通信时间，通信延迟可以被完全隐藏。
但当 $P$ 很大（如 64 块 GPU 处理 128K 序列，每块只有 2K token）时，
计算时间很短，通信延迟无法被隐藏，
Ring Attention 的效率急剧下降。

经验规则是：
\textbf{每块 GPU 至少应分配 $\geq$8K token}
以确保合理的计算--通信比。
这意味着 128K 序列最多使用 16 块 GPU 做 Ring Attention，
1M 序列最多使用 128 块 GPU。
超过这个限制后增加 GPU 的收益很小，
反而因为更多的通信轮次而增加延迟。

\section{长上下文训练的实践}

支持长上下文不是简单地「把序列长度调大」就行。
需要在数据、训练策略和工程三个维度同时发力。

\subsection{多阶段扩展训练}

通常需要分阶段训练：

\textbf{阶段一}：在标准上下文长度（如 4K--8K）上完成主体预训练，
消耗大部分训练 token（如 14T 中的 13.5T）。
这个阶段的重点是学习语言和知识，上下文长度不是关注点。

\textbf{阶段二}：逐步扩展上下文长度，
先扩展到 32K，训练少量数据适应；
然后扩展到 128K，再训练一批数据。
Qwen2.5 的长上下文版本甚至使用了三阶段扩展（32K → 128K → 1M）。
每次扩展时需要调整 RoPE 的频率参数（通过 YaRN 或类似方法）。

LLaMA 3~\cite{dubey2024llama3} 的长上下文训练提供了一个典型的工程实例：
预训练在 8K 上下文上完成，
然后通过六个子阶段逐步扩展到 128K，
每个子阶段将上下文长度翻倍（8K → 16K → 32K → ... → 128K），
同时逐步增大 RoPE 的 base 频率
（从 $5 \times 10^5$ 逐步增大到 $8 \times 10^6$）。

\textbf{阶段三}：在长上下文数据上做 SFT，
确保模型能够在长文本中遵循指令、
回答关于文档特定部分的问题、
以及保持长距离的一致性。

这种分阶段策略的原因是：
长序列的训练效率远低于短序列，
注意力的 $O(n^2)$ 复杂度意味着
一条 128K 的样本比 32 条 4K 的样本\textbf{慢得多}，
但包含的有效 token 数是一样的。
在短序列上完成大部分学习，
然后在长序列上做少量适应，是效率最优的策略。

\subsection{长上下文训练的内存优化}

长上下文训练不仅计算量更大，
对内存的需求也急剧增长。
128K 上下文的训练需要存储每一层的
激活值（用于反向传播的梯度计算），
激活值的内存占用与序列长度成正比，
128K 上下文的激活值比 4K 上下文的大 32 倍。

\textbf{梯度检查点}（Gradient Checkpointing / Activation Recomputation）
是应对这一挑战的核心技术。
其思想是：不存储所有层的激活值，
而是只保存部分“检查点”层的激活值，
在反向传播到其他层时\textbf{重新计算}它们的激活值。
这将内存占用从 $O(L)$（$L$ 为层数）降低到 $O(\sqrt{L})$，
代价是约 33\% 的额外计算量（需要重新执行一次前向传播）。

对于超长上下文训练（128K+），
即使使用梯度检查点，单块 GPU 的内存仍然不够。
此时需要与序列并行（Sequence Parallelism）组合，
将一条长序列拆分到多块 GPU 上，
每块 GPU 只处理序列的一部分。
Ring Attention 和 Zig-Zag Ring Attention
正是实现序列并行的核心通信协议。

另一种优化是\textbf{选择性激活重算}：
不是对所有层等概率地做检查点，
而是根据每一层的内存占用和计算成本
选择性地决定哪些层保存激活值、
哪些层重新计算。
注意力层（产生大的 $O(n^2)$ 中间激活值）
通常是重新计算的优先对象，
而 FFN 层（激活值较小但计算较重）
则优先保存。
Megatron 和 DeepSpeed 等训练框架
都实现了这种精细化的激活重算策略。

\subsection{训练数据的序列打包}

高效的长上下文训练还需要解决一个数据层面的问题：
\textbf{序列长度的高度不均匀}。

即使是专门收集的长文档数据，
不同文档的长度差异也很大，
有些书只有 50K token，有些有 500K token。
如果每个训练 batch 只包含一条固定长度的序列，
那么短文档会浪费大量 padding 空间，
长文档则可能超过 GPU 内存限制。

\textbf{序列打包}（Sequence Packing / Document Packing）
将多条不同长度的文档打包成一个固定长度的序列，
用特殊的分隔符标记文档边界。
注意力计算需要修改为\textbf{不跨越文档边界}，
即文档 A 中的 token 不应该关注文档 B 中的 token，
即使它们在同一个 packed 序列中。
这通过修改注意力掩码（attention mask）实现，
在文档边界处设置掩码屏障。

打包策略的选择对训练效率有显著影响。
\textbf{贪心打包}（按文档长度降序排列，
依次填充到当前最空的序列中）
可以达到 95\%+ 的空间利用率，
但可能导致某些序列全是短文档、
某些序列全是长文档，
batch 内的计算负载不均衡。
\textbf{随机打包}（随机混合长短文档）
负载更均衡但空间利用率稍低。
在实践中，大多数训练框架使用
贪心打包 + 随机 shuffle 的组合策略。

\subsection{长上下文训练数据}

长上下文训练面临一个独特的数据挑战：
\textbf{大多数高质量文本都不长}。
网页平均长度约 2K--4K token，
新闻文章约 500--2000 token。
真正超过 32K token 的高质量文档并不多。

长上下文训练数据的来源：
\begin{itemize}[noitemsep]
  \item \textbf{书籍}：单本书通常在 50K--500K token，
        是最天然的长上下文数据源。
        但版权问题限制了可用数量。
  \item \textbf{代码仓库}：一个完整的代码仓库
        （包含所有文件及其依赖关系）
        可以轻松超过 100K token。
        代码的优势是结构清晰、
        有明确的跨文件引用关系。
  \item \textbf{学术论文}：单篇论文 8K--30K token，
        可以将相关论文串联起来作为长文档。
  \item \textbf{法律和财务文档}：合同、招股书、法规
        动辄数十万 token，是天然的长文本。
  \item \textbf{合成长文档}：将多个相关短文档
        按逻辑顺序拼接成长文档。
        例如将关于同一主题的 Wikipedia 页面拼接起来，
        或者在一组文档上生成跨文档的问答对。
        这种方法需要小心处理拼接边界，
        避免模型学到「文档之间无关联」的错误模式。
  \item \textbf{多轮对话记录}：真实的客服对话、
        技术支持会话、多轮编程协助记录。
        这些数据天然包含长距离的上下文依赖
        （如对话早期提到的约束在后期需要被遵守），
        是长上下文 SFT 的优质数据源。
\end{itemize}

长上下文训练的数据配比是一个容易被忽视但极其重要的问题。
理想的配比需要在几个互相矛盾的目标之间平衡：

\textbf{长度分布的多样性}。
如果训练数据全部是 128K 长度的文档，
模型可能会“过拟合”到 128K 的长度模式，
在短序列上表现反而退化
（因为它期望输入总是很长的，
面对短输入时注意力模式不适应）。
经验法则是在长上下文训练阶段
保持短文档（4K--8K）、中等文档（32K--64K）
和长文档（128K+）各占约 1/3 的比例。

\textbf{内容类型的平衡}。
代码仓库是最容易获取的长上下文数据，
但如果训练数据中代码占比过高，
模型在自然语言长文档理解上可能变差。
相反，自然语言长文档（如书籍、法律文件）
通常质量更高但数量有限。
LLaMA 3~\cite{dubey2024llama3} 的方案是
在长上下文阶段将代码比例从预训练的 $\sim$15\%
提升到 $\sim$30\%（因为代码的跨文件引用
是天然的长距离依赖），
同时增加学术论文和技术文档的比例。

\textbf{有效长程依赖的密度}。
一条 128K token 的文档如果只是简单地将
无关的短段落拼接在一起，
那么它虽然“长”，但不包含真正的长程依赖，
模型不需要关注远处的信息也能做好下一个 token 预测。
这种“虚假长文档”对长上下文能力的训练没有帮助。
高质量的长上下文训练数据应该包含丰富的
\textbf{跨距离引用}，
如小说中的情节回呼、
代码仓库中的跨文件函数调用、
学术论文中的前后文交叉引用。
一些研究团队使用启发式规则
来估计文档中长程依赖的密度，
优先选择依赖密度高的文档进入训练集。

\subsection{长上下文 SFT 数据}

阶段三的长上下文 SFT 需要专门设计的指令数据。
典型的任务类型包括：
\begin{itemize}[noitemsep]
  \item \textbf{长文档总结}：给出一篇 50K+ token 的文档，
        要求模型生成不同粒度的摘要（一句话/一段/多页）。
  \item \textbf{多文档问答}：给出多篇文档，
        要求模型综合多个来源回答问题。
        这测试的是跨文档信息整合能力。
  \item \textbf{仓库级代码理解}：给出整个代码仓库，
        要求模型解释架构、定位 bug、
        或在保持全局一致性的前提下修改代码。
  \item \textbf{长距离事实一致性}：
        在对话的早期提供一组事实，
        在很久之后（上万 token 的对话之后）
        询问与这些事实相关的问题。
\end{itemize}

\subsection{推理工程优化}

长上下文推理的主要瓶颈在 prefill 阶段，
即模型首次处理整个输入序列、计算所有 token 的 KV cache 的阶段。
对于 1M token 的输入，prefill 需要处理的注意力计算量
是 4K 输入的约 $62{,}500\times$（$(10^6/4000)^2$），
即使有 FlashAttention 优化，
朴素实现在 8$\times$H100 上也需要接近一分钟。

SGLang 在 2026 年初发布的 \textbf{Chunked Pipeline Parallelism}
是解决这一瓶颈的重要工程突破。
核心思想是将长序列的 prefill
分成多个 chunk（例如每个 chunk 32K token），
将不同 chunk 分配到不同的 GPU 上，
以流水线方式并行计算。
当 GPU~1 在计算 chunk~1 的注意力时，
GPU~2 可以同时计算 chunk~2 中
不依赖 chunk~1 结果的部分计算
（如 Q/K/V 的线性投影）。
不同 chunk 之间的因果注意力依赖
通过在 GPU 之间传递 KV cache 的中间结果来满足，
这与 Ring Attention 的环形传递机制类似，
但针对 prefill 阶段做了专门优化。

在 DeepSeek-V3.1 上的实测显示，
处理 1M token 的输入延迟从 55.5 秒降至 10.5 秒，
降低了 81\%，相当于 $3.31\times$ 的吞吐提升。
这意味着用户可以在约 10 秒内完成对一整本书的「阅读」，
极大地提升了长上下文模型的实用性。
该优化已集成到 SGLang v0.4+ 中，
与 vLLM 的 chunked prefill 方案互补，
vLLM 侧重于在单 GPU 上减少 prefill 的内存峰值，
SGLang 则侧重于跨 GPU 的流水线并行加速。
两种方案可以组合使用：
先用 SGLang 的流水线将 prefill 分散到多 GPU，
再在每块 GPU 上用 vLLM 风格的分块策略
控制单 GPU 内的内存压力。

\subsection{前缀缓存：跨请求复用}

长上下文推理中另一个关键优化是\textbf{前缀缓存}（Prefix Caching）。
当多个请求共享相同的前缀时
（如同一个系统提示词，或同一份文档被多个用户查询），
可以将前缀的 KV cache 缓存起来，
后续请求直接复用而无需重新计算。

这在长上下文场景中尤其有价值。
想象一个文档分析场景：
用户上传了一份 100K token 的法律合同，
然后连续问了 20 个不同的问题。
如果没有前缀缓存，
每个问题都需要重新处理 100K token 的 prefill，
总计 2M token 的 prefill 计算。
有了前缀缓存，100K token 的 prefill 只需做一次，
后续 19 个问题直接复用 KV cache，
总计算量减少约 20 倍。

vLLM 的 PagedAttention 机制~\cite{kwon2023vllm}% NEW REF: kwon2023vllm = Kwon et al., "Efficient Memory Management for Large Language Model Serving with PagedAttention", SOSP 2023
为前缀缓存提供了高效的内存管理。
KV cache 被组织为固定大小的“页”（如每页 16 个 token），
共享前缀的不同请求引用相同的页，
只有请求特有的后缀部分需要新的内存页。
这种“写时复制”（copy-on-write）的策略
在多用户场景中极大地降低了内存开销。

实测数据表明，在典型的长文档问答场景中，
前缀缓存可以将总推理成本降低 60\%--85\%，
这使得长上下文 API 的定价得以大幅下降，
推动了长上下文应用的商业可行性。

\subsection{推理时的长度感知调度}

长上下文推理的另一个工程挑战是
\textbf{请求长度的高度异质性}。
在一个典型的推理服务中，
不同用户的请求长度可能从 100 token 到 1M token 不等。
如果将长短请求混合在同一个 batch 中，
短请求必须等待长请求完成 prefill 才能开始处理，
这种“队首阻塞”（head-of-line blocking）
严重影响了短请求的延迟。

\textbf{长度感知调度}（Length-Aware Scheduling）
将请求按长度分组处理：
短请求（$<$4K token）在一组 GPU 上快速处理，
中等请求（4K--128K）在另一组 GPU 上处理，
超长请求（$>$128K）单独处理。
这种分流策略确保了
短请求不会被长请求阻塞，
同时长请求可以获得足够的 GPU 资源
来完成 prefill 和 KV cache 管理。

SGLang 和 vLLM 在 2026 年初都引入了
这种长度感知调度策略，
实测将 99 分位（P99）延迟降低了约 40\%，
短请求不再被不可预测的长请求所阻塞。

\section{长上下文的实际应用场景}

长上下文能力开启了许多此前不可能的应用场景：

\textbf{整仓库代码理解}：
将一个完整的代码仓库（数万行代码）放入上下文，
让模型理解整体架构后再回答具体问题或做修改。
这比传统的 RAG（只检索相关片段）
能获得更全面的上下文理解。
Cursor 和 GitHub Copilot Workspace 等开发工具
已经在利用这种能力，
将整个项目的代码、配置文件、README
一起放入上下文，让模型做出全局一致的修改。

\textbf{长文档分析}：
将一份数百页的法律合同或学术论文完整放入上下文，
让模型进行总结、比较或提取特定信息。
这避免了将文档切分成小块后丢失跨段落关联的问题。
在金融领域，分析师可以将一整份招股书
（通常 200--500 页）放入上下文，
直接询问「这家公司的主要风险因素和竞争优势分别是什么？」

\textbf{多文档对比分析}：
将多篇研究论文同时放入上下文，
让模型比较它们的方法、结果和结论。
这在文献综述和系统性综述中特别有价值，
传统做法需要人类研究者反复翻阅多篇论文，
而长上下文模型可以「同时阅读」所有论文并做对比。

\textbf{多轮对话的长期记忆}：
在一次持续数小时的编程对话中，
模型可以记住早期讨论的设计决策，
在后期做修改时保持一致性。
对于短上下文模型，长对话必须依赖
RAG 或摘要来维持「记忆」，
但这些方法不可避免地会丢失信息。
足够长的上下文窗口可以直接保留完整的对话历史。
但这也引入了一个新的挑战：
随着对话历史的增长，
推理成本线性增加（每次生成新 token 都需要
对整个历史做注意力计算），
达到一定长度后实时响应变得困难。
在实践中，通常需要在“保留完整历史”
和“保持响应速度”之间做权衡，
例如将超过一定长度的历史
压缩为摘要或移入长期记忆系统。

\textbf{少样本学习}：
在上下文中放入大量的示例和说明，
让模型在不做任何微调的情况下学会新任务。
更长的上下文意味着可以提供更多、更多样的示例。
一个极端的例子：在上下文中放入整本编程教材，
然后要求模型按照教材的风格和约定写代码，
这比用几个 few-shot 示例效果好得多。

\textbf{长内容生成}：
生成长篇小说章节、技术文档、剧本等需要
在数万 token 的范围内保持情节、角色、
技术概念一致性的内容。
短上下文模型在长篇生成中常出现前后矛盾，
因为它「忘记了」前面写过什么。

\textbf{Agent 的长时间任务执行}：
2025--2026 年 AI Agent 的兴起
为长上下文创造了一个新的高价值应用场景。
一个执行复杂多步骤任务的 Agent
（如完成一个 SWE-bench 编程任务
或执行一系列网页操作）
需要在上下文中维护完整的任务状态：
已完成的步骤、中间结果、失败的尝试及其教训、
当前的目标和约束。
一个持续数小时的 Agent 会话
可能累积超过 100K token 的上下文，
包括工具调用和返回值、思考过程、
错误堆栈信息等。
长上下文确保 Agent 在任务后期
仍然能“记住”前期发现的关键信息，
避免重复犯相同的错误。

\textbf{视频和多模态理解}：
LLaMA 4 等原生多模态模型
将视频帧转化为 token 序列。
一段 10 分钟的视频以每秒 1 帧、
每帧 256 个 token 的方式编码，
就需要约 153K token。
长上下文使模型能“观看”完整的视频片段，
理解前后的时间关系和情节发展，
而不是只看到几秒钟的片段。
这种能力对视频摘要、视频问答、
视频审核等应用至关重要。

但正如下一节将讨论的，
更长的上下文不等于更好的利用。
在实际应用中，如何组织和筛选放入上下文的信息
（context engineering）与上下文长度本身同样重要，
甚至更重要。

\section{RAG vs 长上下文：并非二选一}
\label{sec:rag-vs-long-context}

2025 年的一个热门辩题是：
“长上下文是否让 RAG 过时了？”
如果模型可以一次性读入 1M token（约 1500 页文档），
为什么还需要先检索再阅读的两阶段 RAG 流程？

答案并非非此即彼。
RAG 和长上下文各有不可替代的优势，
最优方案通常是两者的协同。

\subsection{长上下文的优势}

\textbf{上下文完整性}。
长上下文的最大优势是模型能“看到”所有信息。
在 RAG 中，检索器必须在用户提问之前
决定哪些文档片段是相关的：
这是一个本质上不完美的预判。
如果相关信息分散在文档的多个部分，
检索器可能只找到其中一部分，
遗漏关键的关联线索。
长上下文则不存在这个问题，
所有信息都在模型的“视野”中。

\textbf{跨段落推理}。
当答案需要综合文档不同部分的信息时，
长上下文天然支持这种跨段落推理。
RAG 返回的孤立片段之间可能缺少过渡上下文，
模型难以理解它们的逻辑关系。

\textbf{简化工程复杂度}。
RAG 需要构建和维护向量数据库、调优检索器、
设计 chunking 策略、处理文档更新的索引同步等。
长上下文方案只需要将文档塞入 prompt，
工程复杂度大幅降低。

\subsection{RAG 的优势}

\textbf{规模不受限}。
即使是 1M token 的上下文也只能容纳约 5--10 本书。
但一个向量数据库可以索引数百万文档、
数十亿 token 的知识库，
这远超任何上下文窗口的容量。
对于企业级的知识管理场景（如全公司的内部文档），
RAG 是唯一可行的选择。

\textbf{成本效率}。
处理 1M token 的输入需要大量的计算和内存。
以当前的 API 定价，处理 1M token 输入的成本
约为 \$5--\$15（取决于模型和供应商）。
而 RAG 只需处理检索到的少量相关片段
（通常 2K--10K token），成本仅为前者的 1\%--5\%。
对于高频查询场景，成本差异是决定性的。

\textbf{时效性}。
RAG 的知识库可以实时更新，
新文档可以在数秒内被索引和检索。
而长上下文方案中的知识必须在每次查询时重新加载，
没有“增量更新”的概念。

\subsection{最优实践：RAG + 长上下文}

2025--2026 年的生产实践收敛到了一种
\textbf{RAG 筛选 + 长上下文理解}的混合策略：

\begin{enumerate}[noitemsep]
  \item \textbf{粗筛}：使用 RAG 从大规模知识库中
        检索 top-20 到 top-50 个相关文档片段。
  \item \textbf{精排}：使用一个 reranker 模型
        对检索结果按相关性重新排序。
  \item \textbf{理解}：将排序后的文档片段
        （通常 20K--100K token）
        一次性放入长上下文模型中，
        让模型综合所有信息给出回答。
\end{enumerate}

这种三阶段流程结合了两者的优势：
RAG 解决了“从海量数据中找到相关信息”的问题，
长上下文解决了“综合多个信息源做推理”的问题。
经验表明，这种混合方案在准确率和成本之间
取得了比纯 RAG 或纯长上下文更好的平衡。

一个关键的设计参数是\textbf{检索数量}（top-$k$）的选择。
$k$ 太小（如 top-3），可能遗漏相关信息。
$k$ 太大（如 top-100），引入大量噪声。
经验法则：$k$ 应使检索到的总 token 数
占模型有效上下文长度的 30\%--50\%，
留出足够空间给系统提示词、
用户查询和模型的推理过程。
对于 128K 上下文的模型，
top-20 到 top-50 个 512-token 的片段
（总计 10K--25K token）通常是最优配置。

\subsection{“长上下文让 RAG 过时了？”的误解}

这个说法过度简化了问题。
更准确的描述是：
\textbf{长上下文降低了简单 RAG 场景的必要性，
但增加了复杂 RAG 场景的能力}。

在“查找一个具体事实”的简单场景中，
如果文档总量不超过 1M token，
确实可以用长上下文直接替代 RAG，
更简单，更可靠。

但在“从百万文档中综合多个来源做推理”的复杂场景中，
长上下文不仅没有让 RAG 过时，
反而\textbf{增强}了 RAG 的能力，
RAG 检索出的 50 个相关片段
在长上下文中可以被一次性综合理解，
而不需要像短上下文时代那样
逐个片段做“增量式”的推理。

\section{Context Rot 与 NIAH：长上下文的质量评估}

在追求更长的上下文窗口时，
我们需要区分两个不同的能力维度：
\textbf{检索}（retrieval）和\textbf{理解}（reasoning）。
2025 年的研究表明，这两个维度之间的差距
远比人们最初预期的更大。

\subsection{Needle in a Haystack 测试}

NIAH（Needle in a Haystack）是最广泛使用的长上下文评估方法。
测试方法很简单：在一段无关的长文本（haystack）中
插入一条关键信息（needle），
然后询问模型关于这条信息的问题。
通过改变 needle 的位置（开头/中间/结尾）
和 haystack 的长度，
可以绘制出模型在不同位置和长度下的检索准确率热力图。

NIAH 测试揭示了几个重要模式：
\begin{itemize}[noitemsep]
  \item 大多数现代模型在 NIAH 上表现优秀，
        即使在 128K 长度下，主流模型的检索准确率也接近 100\%。
  \item 但 NIAH 本质上只测试\textbf{检索}，
        模型能否找到特定信息。
        这是长上下文能力的\textbf{必要条件}，
        但远不是\textbf{充分条件}。
\end{itemize}

一个模型可以完美地通过 NIAH（找到 needle），
但在需要\textbf{综合}长上下文中多处信息做推理时表现糟糕。
比如「根据文档第 3 章和第 7 章的描述，
这两部分的结论是否矛盾？」：
这要求模型不仅找到两处信息，
还要理解并比较它们的含义。

NIAH 测试的方法论本身也值得详细讨论。
标准的 NIAH 测试使用一个固定的 “needle” 句子
（如 “The best thing to do in San Francisco is eat a sandwich and sit in Dolores Park on a sunny day”），
将其插入到不同位置的 Paul Graham 散文中（“haystack”），
然后询问模型 “What is the best thing to do in San Francisco?”。
通过系统性地扫描 needle 位置（从 0\% 到 100\% 的深度）
和 haystack 长度（从 1K 到模型最大上下文），
生成一个 2D 热力图，
横轴是上下文长度，纵轴是 needle 深度位置，
颜色表示检索准确率。

这种测试方法的问题在于它的“人为性”。
真实应用中的“needle”不是一个与 haystack 完全无关的句子，
而是与上下文内容\textbf{语义相关}的信息片段。
Chroma 的 context rot 研究~\cite{chroma2025contextrot}
发现了一个反直觉的结果：
当 haystack 是\textbf{有逻辑结构}的相关文本时，
检索比在随机无关文本中更加困难，
因为语义相关的干扰信息比随机噪声更容易“迷惑”模型。
这意味着标准 NIAH 的结果\textbf{高估}了模型在真实场景中的检索能力。

\subsection{NIAH 的局限与更好的评估}

NIAH 的根本问题是它把长上下文简化成了一个搜索任务。
在真实场景中，用户很少需要「在长文档中找一句话」——
更常见的是需要\textbf{综合}、\textbf{推理}和\textbf{比较}。

更有挑战性的评估任务包括：
\begin{itemize}[noitemsep]
  \item \textbf{Multi-Needle NIAH}：
        在 haystack 中插入\textbf{多个}相关信息片段，
        要求模型综合所有片段才能回答问题。
        这测试的是信息\textbf{聚合}能力，
        而非单点检索
  \item \textbf{RULER}（2024）：
        包含多种长上下文推理任务——
        变量追踪、多跳推理、
        以及需要整合分散在上下文各处的信息的复合任务
  \item \textbf{长文档摘要质量}：
        给模型一篇 100K+ token 的文档，
        评估摘要是否遗漏了中间部分的关键信息。
        这直接测量 lost-in-the-middle 效应
  \item \textbf{BABILong}：
        基于 bAbI 任务改编的长上下文推理基准，
        需要在长文本中追踪实体状态和关系
\end{itemize}

让我们更深入地理解 RULER 和 BABILong 的设计思路，
因为它们代表了长上下文评估从“检索”走向“推理”的关键转变。

\textbf{RULER}~\cite{hsieh2024ruler}（2024）% NEW REF: hsieh2024ruler = Hsieh et al., "RULER: What's the Real Context Size of Your Language Models?", arXiv:2404.06654, 2024
设计了四类任务来系统性地评估长上下文能力：
（1）\textbf{单点检索}（类似 NIAH）——
找到上下文中的一条特定信息；
（2）\textbf{多点聚合}——
从上下文的多个位置收集同类信息并汇总，
如“统计文中出现的所有城市名”；
（3）\textbf{变量追踪}——
追踪一个变量在长文本中的值变化，
如“x 最初赋值为 5，后来被改为 7，再被改为 3，
最终 x 的值是多少？”；
（4）\textbf{多跳推理}——
需要从上下文的不同部分提取信息，
然后通过逻辑推理得出结论。

RULER 在不同长度上评估模型的一个关键发现是：
\textbf{大多数模型在单点检索上表现优秀
（即使在 128K 长度下准确率也接近 100\%），
但在多跳推理上随长度增加急剧退化}。
例如，GPT-4-128K 在 4K 上下文的多跳推理准确率约 85\%，
但在 128K 上下文下降到约 55\%——
性能下降了 35 个百分点。
这个落差揭示了检索能力和推理能力之间的鸿沟：
模型可以在长文本中“找到”信息，
但将分散的信息“串联起来做推理”的能力
随上下文增长快速衰退。

\textbf{BABILong} 则采用了一种更可控的评估方法。
它基于 bAbI 基准测试的推理任务
（如“Alice 把球放到盒子里，Bob 把球拿到花园，
球现在在哪里？”），
将关键的推理链条嵌入到长度可控的背景文本中。
通过精确控制推理链的长度和“分散度”
（关键信息片段之间隔了多远），
BABILong 可以精确测量模型在不同“推理跨度”下的性能。
其核心发现与 RULER 一致：
简单的 1-hop 推理在 100K+ 上下文下仍然可靠，
但 3-hop 以上的推理在 32K 以上就开始显著退化。

\subsection{Context Rot：系统性退化}

2025 年 7 月，Chroma 发布了影响深远的
\textbf{context rot} 研究报告~\cite{chroma2025contextrot}，
对 18 个 frontier 模型（包括 GPT-4.1、Claude 4、
Gemini 2.5、Qwen3 等当时最新的模型）
进行了系统性的长上下文评估。
核心发现令人警醒：
\textbf{所有模型都随着输入长度的增加而性能下降，
即使远未达到标称的上下文窗口上限}。

具体发现：
\begin{itemize}[noitemsep]
  \item \textbf{U 型曲线持续存在}：
        模型对序列开头和结尾的信息关注度最高，
        中间部分的信息检索准确度下降超过 30\%。
        这被称为「lost in the middle」效应~\cite{liu2023lost}。
        尽管 2024--2025 年的训练改进缓解了部分问题，
        但 U 型模式并未消失——只是变得不那么陡峭
  \item \textbf{反直觉发现}：有逻辑结构的上下文
        比随机打乱的上下文\textbf{更难}检索——
        模型在有意义的干扰信息中更容易「迷路」。
        这意味着真实应用场景
        （放入完整文档而非随机文本）的难度
        比 NIAH 测试暗示的更高
  \item \textbf{推理 vs 检索的鸿沟}：
        在需要跨上下文推理的任务中，
        性能下降比纯检索任务更为严重。
        模型可能能找到两个相关段落，
        但无法正确地将它们联系起来做推理。
        对于简单的重复词计数任务，
        一些模型在输入增长到几千 token 时
        就开始出现可靠性问题
  \item \textbf{标称 vs 有效上下文}：
        实际测试显示，
        许多模型的「有效上下文利用率」（实际可靠使用的比例）
        仅为标称上下文窗口的 50--80\%。
        例如一个标称 200K 的模型，
        在约 100--160K token 后可靠性就开始下降
  \item \textbf{核心结论}：
        \textbf{Context engineering（如何组织上下文）
        比 context length（能塞多少）更重要}
\end{itemize}

Context rot 研究的一个重要方法论贡献是
它提出了\textbf{有效上下文长度}（effective context length）
的量化方法。
不同于标称上下文窗口（由位置编码和 KV cache 容量决定），
有效上下文长度定义为：
模型在该长度下的任务性能
不低于短上下文基线的 90\%。
用这个标准衡量，
许多标称 128K 的模型的有效上下文
仅为 64K--100K，
标称 1M 的模型有效上下文可能只有 200K--500K。

这个指标为模型选型提供了实际参考：
用户应该关注\textbf{有效上下文长度}而非\textbf{标称上下文窗口}。
一个有效上下文 200K 的模型
在处理 200K 以内的输入时
通常比一个标称 1M 但有效只有 150K 的模型更可靠。

Chroma 的研究还发现了一个有趣的模型间差异：
不同模型的 context rot 曲线形状不同。
有些模型（如 Claude 4 系列）展现出“平坦后陡降”的模式——
在有效上下文范围内性能几乎不下降，
但一旦超过某个临界长度，性能急剧恶化。
另一些模型（如早期的 GPT-4 128K）展现出“渐进退化”模式——
性能从一开始就随长度缓慢下降，
没有明显的“悬崖”。
前者可能更适合“全有或全无”的场景
（要么整个文档放得下，要么不用），
后者可能更适合“尽量多放”的场景
（即使性能有些下降也比不放好）。

\subsection{2026 长上下文实测：MRCR v2 + RULER + LongBench v2 三件套}
\label{sec:long-context-benchmarks-2026}

到 2026 年，原始 NIAH 已被前沿模型刷到饱和，
事实上失去区分能力。
新的"三件套"取代了 NIAH 作为长上下文评估的工业基准：
\textbf{RULER}（结构化检索 + 多跳推理 + 多点聚合 + 变量追踪）、
\textbf{LongBench v2}（真实长推理任务，含中英文）、
\textbf{MRCR v2}（Multi-Round Co-Reference Resolution，
要求在 1M 上下文中跨多轮对话追踪指代）。

\textbf{RULER}~\cite{hsieh2024ruler} 在
llm-stats RULER 1M 排行榜由
\textbf{MiniCPM-SALA}（9B）以 \textbf{0.863} 领跑，
小模型 + 精巧训练再次证明可与百亿级模型在长上下文上正面竞争。
大多数模型在 RULER 上 32K 处即跌破
Llama-2-7B 4K 的 85.6\% 阈值，
意味着\textbf{大多数标称 128K--200K 的模型，
在 32K 之后就失去了基础检索质量}。

\textbf{LongBench v2 长推理}方面，
Gemini 3 Pro 达到 68.2\%、GPT-5.2 达到 54.5\%。
arXiv:2604.07981~\cite{longbenchv2decomp2026} 的 decomposition 方法
把 Loogle、Loong、LongBench-v2、BrowsecompLong、Ruler-qa2、MRCR
六个基准的平均分从 46.3\% 提升到 54.0\%（+7.7 pp），
说明\textbf{显式分解长推理}本身就是一个有效的元策略。

\textbf{MRCR v2 8-needle 1M} 是 2026 年最硬的长上下文基准。
Anthropic 在 2026-03-13 把 1M context GA、
\textbf{取消 long-context premium}~\cite{anthropic2026context1m}：
Opus 价格 \$5/\$25 per M、Sonnet \$3/\$15。
Claude Opus 4.6 在该基准的成绩极其反常：
\textbf{1M MRCR v2 8-needle 取得 76.0--78.3\%}（不同 run 区间），
而它在 128K 上的成绩仅 71.9\%——
\textbf{随上下文增加反而提升}，
这是首次在公开基准上观察到的现象，
来源于 Opus 4.6 的长上下文专项后训练。
对比 Claude Sonnet 4.5 在 1M 仅 18.5\%
（Opus 4.6 ≈ \textbf{4.2$\times$} 提升），
GPT-5.4 从 256K 79.3\% 跌到 1M 36.6\%，
Gemini 3 Pro 1M 仅 24.5\%。
\textbf{256K 4-needle} 这个稍易的赛道上，
GPT-5.2 ≈ \textbf{98\%} 首达近满分，
Opus 4.6 ≈ 91.9--93.0\%。

实测延迟同样揭示了 1M 上下文的真实代价。
tokenmix.ai 2026-04-20 数据~\cite{tokenmix2026latency}：
\textbf{1M tokens prefill latency 已是数十秒级常态}——
Opus 4.6: 60--120 s；Sonnet 4.6: 45--90 s；
Gemini 1.5 Pro: 2 分钟+；Gemini 3.1 Pro: 90--150 s。
500K → 45--90 s；100K → 5--15 s。
GPT-5.4 在 128K（extended 400K）段落保持 \textless 30 s，
是延迟最优的旗舰。
这一组数据意味着：
\textbf{1M 上下文不是"按下回车 1 秒后就能聊天"的体验，
而是"提交后去倒杯咖啡"}——
任何依赖 1M 上下文的产品都需要为长 TTFT 重新设计 UX。

\begin{warnbox}[Needle 已被 game：评估方法的代际更替]
单 needle 的 NIAH 在 2026 年已经失去价值——
所有前沿模型都能在 1M 长度下找到 needle。
但找到 needle $\neq$ 理解上下文：
MRCR v2 多 needle 1M 的成绩从 18.5\% 到 78.3\%
横跨整整 4$\times$，
而这 4 个数都对应"声称支持 1M 上下文"的模型。
评估方法\textbf{必须}从 NIAH 升级到
RULER + LongBench v2 + MRCR v2 三件套，
否则模型选型时会基于失真的"满分错觉"做错决策。
\end{warnbox}

\subsection{Context Engineering：从「更长」到「更智能」}

Anthropic 在 2025 年下半年
明确将 “context engineering” 定义为构建高效 AI Agent 的核心技能，
标志着行业思维从「追求更长的上下文窗口」
转向「更智能地利用有限的上下文预算」~\cite{anthropic2025contexteng}。

\begin{whybox}[Context Engineering 的核心原则]
基于 context rot 的研究和生产实践，
以下原则已成为行业共识：
\begin{itemize}[noitemsep]
  \item \textbf{重要信息前置}：
        把最关键的指令和信息放在上下文的开头，
        利用模型对开头信息的高关注度。
        系统提示词（system prompt）和核心约束应在最前面
  \item \textbf{冗余而非堆积}：
        对关键信息做适度冗余
       （比如在不同位置重复强调重要约束），
        而不是试图塞入尽可能多的不同信息。
        每个不必要的 token 都在\textbf{主动}降低其他 token 的注意力预算
  \item \textbf{结构化组织}：
        使用清晰的标题、分隔符和层次结构
        帮助模型定位和组织信息。
        XML 标签、Markdown 标题、编号列表
        都能显著提高信息检索的可靠性
  \item \textbf{即时上下文（Just-in-Time Context）}：
        不要在一开始就加载所有信息，
        而是在需要时才检索和注入——
        这对 AI Agent 尤为重要。
        一个执行 20 步任务的 Agent，
        在每一步只需要当前相关的上下文，
        而非全部 20 步的信息
  \item \textbf{上下文压缩（Compaction）}：
        对长时间运行的 Agent，
        定期将已处理的上下文压缩为简洁的摘要，
        丢弃不再需要的细节。
        这类似于人类的「工作记忆管理」
  \item \textbf{RAG + 长上下文的协同}：
        RAG 负责从海量文档中筛选相关内容，
        长上下文负责一次性理解所有筛选出的内容。
        两者互补而非替代。
        经验表明 RAG 检索 top-20 到 top-50 相关片段
        并放入 32K--128K 的上下文中，
        效果通常优于直接将整个文档塞入 1M 上下文
\end{itemize}
\end{whybox}

\section{Context Engineering 的信息论框架}
\label{sec:context-engineering-theory}

2026 年 3 月发表的
"The Root Theorem of Context Engineering"
~\cite{contextroottheorem2026}（arXiv:2604.20874）
首次给 context engineering 提供了\textbf{严格的信息论基础}。
论文将 context engineering 形式化为一个根本问题：
\textbf{在有界有损通道内最大化 signal-to-token 比}。
通道有界来自模型的上下文窗口，
通道有损来自注意力分布的稀释（U 型效应、context rot）和
softmax 的归一化竞争。

定理给出了五条推论，构成 2026 年长上下文实践的理论基石：

\begin{itemize}[noitemsep]
  \item \textbf{推论 1（质量单调下降）}：
        质量函数 $F(P)$ 单调随注入 token 数 $|P|$ 下降。
        \emph{推论}：每个不必要的 token 都\textbf{主动}降低其他 token 的注意力预算
  \item \textbf{推论 2（独立优化）}：
        signal（有效信息）与 token 数（容器大小）是\textbf{两个独立的优化变量}。
        加更多 token 不等于加更多 signal，
        signal 密度才是关键
  \item \textbf{推论 3（fidelity gate 必要）}：
        必须有一个由 fidelity 阈值触发的 gate，
        在内容低于阈值时拒绝注入。
        无门 = 信噪比单调下降
  \item \textbf{推论 4（homeostatic persistence）}：
        \textbf{积累-压缩-改写-丢弃}（accumulate-compress-rewrite-discard）
        是\textbf{唯一}能持续维持理解的架构。
        这正是 Anthropic 的 compaction 策略
        和 Agent 记忆系统所做的事
  \item \textbf{推论 5（外部验证必要）}：
        压缩机制运行在它要压缩的通道内部，
        因此\textbf{无法自我验证}——
        必须由外部机制（如独立 evaluator agent）做 fidelity check
\end{itemize}

这五条推论解释了为什么"塞更多上下文"反而更糟、
为什么 Agent 记忆需要主动遗忘、
为什么 long-context inference 必须配合外部验证。
Context engineering 从经验技艺正式升级为
\textbf{有公理基础的工程学科}。

\section{长上下文后训练：OPSDL / StructKV / LongAct}
\label{sec:long-context-posttrain}

2026 年第二季度，长上下文\textbf{后训练}方法集中成熟。
此前的常规做法是：在长样本上做 SFT，
让模型从短上下文模型继续学习处理长输入。
但这种朴素 SFT 面临两个老问题：
（1）长样本数据稀缺、质量参差；
（2）长样本上的 token-level loss 信号被稀释，
关键信息的梯度信号弱。
2026 年三个工作分别从蒸馏、KV cache 结构、注意力 saliency 三个角度
重新设计了长上下文后训练。

\textbf{OPSDL}（On-Policy Self-Distillation for Long Context）
~\cite{opsdl2026}（arXiv:2604.17535）
解决"教师"问题。
长上下文后训练的核心矛盾是：
强短上下文的模型版本（teacher）已存在，
但弱长上下文的同模型版本（student）需要训练；
让一个不擅长长上下文的模型自己监督自己显然不合理。
OPSDL 让\textbf{模型自身的强短上下文能力}充当 teacher：
对每条训练样本，分别用短上下文路径（teacher）
和长上下文路径（student）生成同一个 token 的分布，
通过 \textbf{per-token reverse KL} 监督长上下文 generation
向短上下文 generation 对齐。
"教自己"的关键洞察是：
模型在短上下文上知道正确答案，
长上下文只是把这个知识用起来变难了，
而非答案本身改变。

\textbf{StructKV}~\cite{structkv2026}（arXiv:2604.06746）
从 KV cache 结构入手。
论文识别出 transformer 长序列中的
\textbf{global information hubs}——
少数 token 的 KV 在多层注意力中
被显著高频地引用。
这些"hub token"通常是长文档的标题、章节首段、
关键命名实体等。
StructKV 在 RULER/LongBench 长上下文 SFT 中
显式保留这些 hub token 的 KV，
而对其他 token 做更激进的压缩。
该方法在不增加训练成本的前提下，
让 SFT 后的模型在 hub-driven 检索任务上明显优于均匀处理基线。

\textbf{LongAct}~\cite{longact2026}（arXiv:2604.14922）
从 RL 信号出发。
观察是：长上下文 RL 中，
大多数 token 的梯度贡献接近零，
只有少数"高幅值激活"（high-magnitude activations）
真正驱动学习。
LongAct 用注意力中的高幅值激活作为
\textbf{saliency-guided 稀疏 RL 更新信号}，
仅对这些 token 计算和应用梯度。
LongBench v2 上 LongAct \textbf{+8\%}，
是 2026 年第一个把 RL post-train
应用到长上下文场景并取得显著提升的方法。

三者正交：OPSDL 解决监督信号问题，
StructKV 解决 KV 表征问题，
LongAct 解决 RL 信号问题。
它们共同把长上下文后训练
从"用更多长样本继续训"
推进到"用结构化策略重新设计训练目标"。

\section{稀疏注意力的新进展}
\label{sec:sparse-attention-2026}

2026 年初，多项研究在稀疏注意力方向取得了突破，
为长上下文推理提供了更高效的计算范式。
这些工作的共同目标是：
在不显著损失质量的前提下，
减少长序列中实际参与注意力计算的 token 数量。

图~\ref{fig:attention-patterns} 对比了四种主要的注意力模式。

\begin{figure}[htbp]
\centering
\begin{tikzpicture}[
  >=Stealth,
  cell/.style={minimum size=4mm, inner sep=0pt},
  label node/.style={font=\scriptsize, text=figGray},
  title node/.style={font=\small\bfseries, text=figGray},
]

% Helper macro for drawing attention matrices
% Full Dense Attention
\begin{scope}[shift={(0,0)}]
  \node[title node] at (1.2, 2.8) {Full Dense};
  \node[label node, rotate=90] at (-0.3, 1.2) {Query};
  \node[label node] at (1.2, -0.3) {Key};
  % Draw 6x6 grid, lower triangular filled (causal)
  \foreach \i in {0,...,5} {
    \foreach \j in {0,...,5} {
      \ifnum\j>\i
        \node[cell, fill=figLightGray, draw=white] at (\j*0.4+0.2, 2.0-\i*0.4) {};
      \else
        \node[cell, fill=figBlue!60, draw=white] at (\j*0.4+0.2, 2.0-\i*0.4) {};
      \fi
    }
  }
  \node[label node, text width=22mm, align=center] at (1.2, -0.9) {$O(n^2)$ 计算\\每个 token 看\\所有之前的 token};
\end{scope}

% Sliding Window Attention
\begin{scope}[shift={(3.5,0)}]
  \node[title node] at (1.2, 2.8) {Sliding Window};
  % Draw 6x6 grid, only window=2 around diagonal
  \foreach \i in {0,...,5} {
    \foreach \j in {0,...,5} {
      \pgfmathtruncatemacro{\diff}{\i-\j}
      \ifnum\j>\i
        \node[cell, fill=figLightGray, draw=white] at (\j*0.4+0.2, 2.0-\i*0.4) {};
      \else
        \ifnum\diff>2
          \node[cell, fill=figLightGray, draw=white] at (\j*0.4+0.2, 2.0-\i*0.4) {};
        \else
          \node[cell, fill=figGreen!60, draw=white] at (\j*0.4+0.2, 2.0-\i*0.4) {};
        \fi
      \fi
    }
  }
  \node[label node, text width=22mm, align=center] at (1.2, -0.9) {$O(n \cdot w)$ 计算\\每个 token 只看\\最近 $w$ 个 token};
\end{scope}

% Dilated/Sparse Attention
\begin{scope}[shift={(7.0,0)}]
  \node[title node] at (1.2, 2.8) {Dilated Sparse};
  % Draw 6x6 grid, sparse pattern
  \foreach \i in {0,...,5} {
    \foreach \j in {0,...,5} {
      \pgfmathtruncatemacro{\diff}{\i-\j}
      \pgfmathtruncatemacro{\modval}{mod(\diff,2)}
      \ifnum\j>\i
        \node[cell, fill=figLightGray, draw=white] at (\j*0.4+0.2, 2.0-\i*0.4) {};
      \else
        \ifnum\diff<2
          \node[cell, fill=figOrange!60, draw=white] at (\j*0.4+0.2, 2.0-\i*0.4) {};
        \else
          \ifnum\modval=0
            \node[cell, fill=figOrange!60, draw=white] at (\j*0.4+0.2, 2.0-\i*0.4) {};
          \else
            \node[cell, fill=figLightGray, draw=white] at (\j*0.4+0.2, 2.0-\i*0.4) {};
          \fi
        \fi
      \fi
    }
  }
  \node[label node, text width=22mm, align=center] at (1.2, -0.9) {$O(n \cdot \sqrt{n})$ 计算\\局部密集 +\\远距离稀疏采样};
\end{scope}

% Hierarchical (Double-P style)
\begin{scope}[shift={(10.5,0)}]
  \node[title node] at (1.2, 2.8) {Hierarchical};
  % Draw 6x6 grid, block-sparse pattern
  \foreach \i in {0,...,5} {
    \foreach \j in {0,...,5} {
      \ifnum\j>\i
        \node[cell, fill=figLightGray, draw=white] at (\j*0.4+0.2, 2.0-\i*0.4) {};
      \else
        % Local block (last 2) always attended
        \pgfmathtruncatemacro{\diff}{\i-\j}
        \ifnum\diff<2
          \node[cell, fill=figRed!60, draw=white] at (\j*0.4+0.2, 2.0-\i*0.4) {};
        \else
          % First block always attended (attention sink)
          \ifnum\j<2
            \node[cell, fill=figRed!40, draw=white] at (\j*0.4+0.2, 2.0-\i*0.4) {};
          \else
            \node[cell, fill=figLightGray, draw=white] at (\j*0.4+0.2, 2.0-\i*0.4) {};
          \fi
        \fi
      \fi
    }
  }
  \node[label node, text width=22mm, align=center] at (1.2, -0.9) {自适应计算\\粗筛选块级区域\\精排选 token 级};
\end{scope}

\end{tikzpicture}
\caption{四种注意力模式的对比。颜色区域表示实际参与注意力计算的 token 对。
Full Dense 计算所有因果位置对（$O(n^2)$）。
Sliding Window 只看最近 $w$ 个 token（$O(nw)$）。
Dilated Sparse 在局部保持密集、远距离稀疏采样。
Hierarchical（如 Double-P）先用块级筛选定位相关区域，
再在区域内做精细计算。}
\label{fig:attention-patterns}
\end{figure}

\subsection{Token Sparse Attention}

Token Sparse Attention~\cite{tokensparse2026}（arXiv:2602.03216）
提出了一种\textbf{逐头 Q/K/V 压缩}策略。
不同于 NSA 在 token 级别做选择，
Token Sparse Attention 在每个注意力头内部
对 Q、K、V 向量进行独立的压缩，
根据每个头学到的重要性分数
动态决定哪些 token 的表示需要保留完整精度，
哪些可以被压缩或跳过。
这种逐头的精细化控制
避免了“一刀切”式的全局稀疏策略
在某些头上过度丢弃信息的问题。

这种设计的动机源自一个关键观察：
Transformer 中不同注意力头的功能高度异质化。
部分头专注于局部语法模式（相邻 token 之间的关系），
部分头负责长距离信息检索（如指代消解），
还有部分头在大多数输入上几乎不激活
（“休眠头”现象）。
如果对所有头施加相同的稀疏策略，
检索头可能丢失关键的远距离信号，
而休眠头上的保留操作则纯属浪费。
Token Sparse Attention 通过为每个头学习独立的重要性评估函数
$\sigma_h: \mathbb{R}^{d_h} \to [0, 1]$（$h$ 为头索引），
让模型自主决定每个头在每个位置的压缩力度。

具体实现上，重要性分数的计算引入的额外开销极小——
每个头仅需一个轻量的线性投影加 sigmoid 激活，
参数量不到原始注意力权重的 0.1\%。
在 128K 序列长度的基准测试中，
Token Sparse Attention 相比标准全注意力
减少了约 45\%--60\% 的实际注意力计算量，
同时在困惑度（perplexity）上的退化不超过 0.3\%。
值得注意的是，这种方法可以与 FlashAttention
的分块计算策略兼容——
被跳过的 token 不参与分块内的 SRAM 读写，
从而在 IO 层面也获得了加速。

\subsection{Double-P 分层稀疏注意力}

Double-P~\cite{doublep2026}（arXiv:2602.05191）
引入了\textbf{分层 top-$p$ 稀疏注意力}机制。
其核心思想是将稀疏选择分为两个层次：
第一层 top-$p$ 在粗粒度（token 块级别）上筛选相关区域，
第二层 top-$p$ 在精细粒度（单 token 级别）上
从候选区域中选择最相关的 token。
这种两级筛选策略
既保持了全局上下文的感知能力，
又避免了在完整序列上做精细排序的计算开销。

Double-P 的名字来源于两次 top-$p$ 采样——
与常见的 top-$k$（保留固定数量的元素）不同，
top-$p$ 保留\textbf{累积概率达到阈值 $p$} 的元素集合。
这意味着当注意力分布集中时（少数 token 高度相关），
实际参与计算的 token 数量很少；
当注意力分布分散时（信息分布在多个位置），
保留的 token 数量自动增加。
这种自适应的稀疏度控制
比固定 top-$k$ 更适合真实文本中高度变化的注意力模式。

第一层的粗粒度筛选将序列划分为固定大小的块
（通常为 128 或 256 个 token），
对每个块计算一个聚合的相关性分数
（块内所有 key 向量的均值与 query 的点积），
然后通过 top-$p_1$ 筛选出候选块。
第二层在候选块内部执行精细的 token 级 top-$p_2$ 筛选，
最终只对选中的 token 计算完整的注意力。
两层筛选的总计算复杂度为
$O(n/B + B \cdot |\mathcal{C}|)$，
其中 $B$ 是块大小，$|\mathcal{C}|$ 是候选块数量——
在典型配置下，这远低于全注意力的 $O(n)$。

实验结果表明，Double-P 在 256K 上下文长度上
相比全注意力减少了约 70\% 的计算量，
而在 NIAH 和 Multi-Needle 检索任务上
性能下降不超过 1\%。
相比 Token Sparse Attention 的逐头策略，
Double-P 的优势在于其分层结构更适合超长序列——
粗粒度筛选在序列很长时能快速排除大部分不相关区域，
避免了在百万级 token 上做逐个排序的高延迟。

\subsection{DySCO：动态注意力缩放}

DySCO（Dynamic Scaling for Context Optimization）
~\cite{dysco2026}（arXiv:2602.22175）
提出了一种基于\textbf{检索头}（retrieval heads）
的动态注意力缩放机制。
研究发现，Transformer 中存在少数“检索头”——
这些头在长上下文任务中负责精确检索远距离信息。
DySCO 识别这些检索头后，
动态调整不同头的注意力计算范围：
检索头保持对完整上下文的注意力，
非检索头使用更激进的稀疏策略。
这种自适应策略在不同任务上
自动找到效率和质量的平衡点。

DySCO 的“检索头”概念建立在此前的实证研究之上。
多项工作已经观察到，
在一个典型的 32 头 Transformer 模型中，
仅有约 3--5 个头在长距离检索任务中表现出
显著高于随机基线的检索能力。
这些检索头具有独特的注意力分布特征——
它们的注意力权重分布更加“尖锐”（低熵），
倾向于将大部分注意力集中在少数关键位置上，
而非像其他头那样广泛分散。
DySCO 利用这一特征作为检索头的识别信号：
在推理时，通过监测每个头的注意力熵，
实时判断哪些头正在执行检索操作。

DySCO 的缩放机制在实现上采用了两级控制。
\textbf{空间缩放}决定每个头的注意力窗口大小——
检索头使用完整的上下文窗口（如 128K），
模式匹配头使用中等窗口（如 8K--16K），
局部语法头使用短窗口（如 1K--2K）。
\textbf{精度缩放}决定注意力计算的数值精度——
检索头使用 FP16/BF16 保证精确的分数排序，
非检索头可以使用更低精度（如 FP8）
以进一步节省计算和内存。
两级缩放的组合使得 DySCO 在 128K 上下文上
相比全注意力实现了约 $2.1\times$ 的端到端加速，
同时在 RULER 和 BABILong 等综合评估中
保持了 98\% 以上的原始性能。

与 Token Sparse Attention 和 Double-P 的比较揭示了
三种方法的互补性：
Token Sparse Attention 关注“每个头内部哪些 token 重要”，
Double-P 关注“全局哪些区域包含相关信息”，
DySCO 关注“哪些头需要看到完整上下文”。
原则上，这三种策略可以正交地组合——
先用 DySCO 确定每个头的注意力范围，
再在检索头内用 Double-P 做分层筛选，
在其他头内用 Token Sparse 做逐头压缩。
这种组合策略是稀疏注意力研究的一个开放方向。

\subsection{“Lost in the Middle” 的精确理论}

2026 年 3 月，“Lost in the Middle”
效应首次获得了严格的理论解释~\cite{lostmiddle2026}（arXiv:2603.10123）。
研究者从 Transformer 的因果注意力结构出发，
数学地证明了 U 型注意力分布是
自回归训练目标和 softmax 归一化
在有限精度下的\textbf{必然结果}——
而非模型容量不足或训练数据不足导致的偶然现象。

论文的核心推导可以直觉地理解为：
在因果注意力中，位置 $t$ 的 token
只能关注位置 $1$ 到 $t$ 的 token。
这意味着位置 1 的 token 出现在\textbf{所有}后续 token 的注意力计算中——
它被“参考”了 $n-1$ 次（$n$ 为序列长度）。
而位置 $n/2$（中间）的 token 只被后半部分的 token 参考，
参考次数约为 $n/2$。
训练过程中，梯度信号的累积
使得开头位置的信息被反复强化。
同时，softmax 的归一化特性
意味着随着参与竞争的 token 数量增加，
每个 token 分到的注意力“预算”被稀释——
但结尾位置因近因效应（recency bias）获得额外优势。
中间位置既没有累积强化的优势，
也没有近因效应的加持，
因此系统性地被低估。

论文进一步证明了 U 型效应的\textbf{不可避免性}：
只要同时满足三个条件——
（1）因果注意力掩码、
（2）softmax 归一化、
（3）有限的注意力头维度——
U 型分布就是信息论意义上的必然结果。
这排除了通过简单调整超参数
（如学习率、训练步数、数据分布）来消除该效应的可能性。

这一理论结果意味着：
在不改变基本架构的前提下，
仅靠增加训练数据或模型参数
\textbf{无法彻底消除} lost-in-the-middle 效应。
解决方案必须来自架构层面的创新——
如 iRoPE 的 NoPE 层（不使用因果掩码中的位置偏置）、
稀疏注意力（改变注意力的竞争结构）、
分层处理（在不同粒度上分别计算注意力）、
或下文将讨论的长期记忆机制
（将知识检索从注意力计算中分离出来）。
该论文也为 context engineering 的“重要信息前置”原则
提供了理论基础——
既然 U 型效应不可消除，
那么在实践中就应该主动利用它，
将关键信息放在模型天然高关注的位置。

\subsection{UltraLong-8B：从 128K 到 4M 的高效训练}

UltraLong-8B~\cite{ultralong2025}（arXiv:2504.06214）
证明了一个令人振奋的可能性：
通过精心设计的多阶段训练策略，
一个仅有 8B 参数的模型可以从 128K
高效地扩展到 \textbf{4M token} 的上下文长度。
关键技术包括渐进式 RoPE 频率调整、
分阶段长度扩展训练、以及
针对超长序列的数据采样策略。
这表明超长上下文并不必然需要超大参数模型——
训练方法论的创新可以让中等规模模型
也具备百万级上下文处理能力。

UltraLong-8B 的训练流程分为四个精心设计的阶段，
每个阶段都有明确的目标和优化策略。
\textbf{阶段一}（128K $\to$ 512K）使用温和的 RoPE base 频率调整
（从 $5 \times 10^5$ 增加到 $2 \times 10^6$），
在大约 2B token 的长文档数据上继续训练。
\textbf{阶段二}（512K $\to$ 1M）进一步增大 base 频率
并引入 YaRN 的分频段插值策略，
训练约 1B token。
\textbf{阶段三}（1M $\to$ 2M）减小学习率并增加序列打包密度，
确保每个 batch 中包含足够的有效长程依赖。
\textbf{阶段四}（2M $\to$ 4M）采用极保守的训练配置
——极低学习率、大梯度累积步数——
在约 500M token 上完成最终适应。
整个扩展过程的训练成本不到原始预训练的 2\%。

这项工作的核心洞察是
\textbf{长度扩展的难度并非均匀增长}。
从 128K 扩展到 512K 相对容易——
模型在预训练中已经见过大量中等长度的模式，
只需微调位置编码的范围。
但从 1M 到 4M 则困难得多，
因为模型需要学习此前从未遇到的超远距离依赖关系。
UltraLong-8B 通过在每个阶段精确控制
学习率衰减曲线和数据配比
（长文档与短文档的混合比例），
避免了“灾难性遗忘”——
即模型在学会处理 4M 上下文的同时
丧失了短上下文任务的性能。

在评估方面，UltraLong-8B 在 4M 上下文长度下
NIAH 检索准确率仍保持 95\% 以上，
在 RULER 多跳推理任务上
相比 128K 基线模型仅下降约 8\%。
考虑到模型只有 8B 参数——
远小于 LLaMA 4 Scout（109B）或 DeepSeek-V3（671B）——
这一结果具有重要的实践意义：
它为资源受限的研究团队
提供了一条可行的长上下文探索路径，
证明了“小模型 + 精巧训练”
在长上下文维度上可以与大模型竞争。

\section{长期记忆：超越传统上下文窗口}
\label{sec:long-term-memory}

2026 年，长上下文领域迎来了一个根本性的范式转变——
从“如何让上下文窗口更长”
转向“如何让模型具有真正的长期记忆”。
传统的上下文窗口无论多长，
都受限于一个根本约束：
窗口内的所有 token 必须参与注意力计算，
计算和内存成本与窗口大小成正比（甚至平方关系）。

长期记忆（Long-Term Memory, LTM）
的核心思想是将“记忆”与“计算”分离——
让模型拥有一个远大于上下文窗口的
持久化知识存储，
在需要时以极低的成本检索。

\subsection{DeepSeek Engram：条件记忆}

2026 年 1 月 12 日，DeepSeek 发表了
一篇可能重塑 LLM 架构范式的论文——
“Conditional Memory via Scalable Lookup:
A New Axis of Sparsity for LLMs”
~\cite{engram2026}（arXiv:2601.07372），
同样由 CEO 梁文锋联合署名。

\subsubsection{核心洞察：Transformer 缺少查找原语}

Engram 论文的出发点是一个深刻的观察：
\textbf{Transformer 架构缺少原生的知识查找原语}。
当前模型回忆事实知识时，
必须通过多层注意力和 FFN 的计算来“模拟”检索——
这是一种极其低效的方式。
就好比一个图书馆员
每次找书都不查目录，
而是把所有书从头到尾翻一遍。

Engram 模块通过引入一种
\textbf{现代化的 N-gram 嵌入}实现了 $O(1)$ 的查找操作。
给定输入的局部上下文（几个相邻 token），
Engram 直接从一个巨大的嵌入表中
查找对应的知识表示——
不需要经过注意力计算。

为什么“查找”比“计算”更高效？
考虑一个简单的例子：
当模型遇到 “The capital of France is” 这个序列时，
传统 Transformer 需要通过多层注意力和 FFN 的计算
来“回忆”出答案 “Paris”——
这需要在权重矩阵中执行多次矩阵乘法。
而 Engram 可以直接用 “capital of France” 这个 N-gram
在嵌入表中查找对应的向量表示——
只需一次哈希查找和向量读取，
计算量与知识库大小无关。

Engram 的“N-gram”并非传统的固定长度 N-gram。
它使用了一种可学习的上下文窗口
（论文称之为 “conditional context”），
模型自主决定需要多长的局部上下文
来唯一标识当前需要查找的知识。
对于常见的、简单的知识
（如 “capital of France $\to$ Paris”），
2--3 个 token 的上下文就足以唯一标识。
对于罕见的、复杂的知识
（如某个特定论文中的特定实验结果），
可能需要 5--8 个 token 的上下文。
这种自适应性使得 Engram 在不同知识复杂度上
都能保持高效。

\subsubsection{U 型缩放定律}

论文最引人注目的发现之一是
关于稀疏模型容量分配的\textbf{U 型缩放定律}：
在给定的总参数预算下，
分配给 Engram 记忆的最优比例约为 \textbf{75\%}——
也就是说，一个高效的稀疏模型
应该将四分之三的容量用于存储可查找的知识，
仅用四分之一的容量用于计算。

这个发现挑战了现有模型的设计惯例——
当前 LLM 几乎将所有参数都用于计算
（注意力和 FFN 层），
没有专门的知识存储空间。

\subsubsection{部署优势}

Engram 的另一个重要特性是
其存储可以放在\textbf{主机内存}（host memory）上
而非昂贵的 HBM 上。
即使是 100B 参数规模的 Engram 存储，
放在主机内存中也仅带来\textbf{不到 3\% 的吞吐量损失}——
因为 Engram 的查找操作是简单的表查询，
不需要 GPU 的高速计算能力。
这意味着在相同的 GPU 硬件上，
可以“免费”地为模型添加
数十倍于 GPU 显存容量的知识存储。

DeepSeek 已在 GitHub 上以 Apache 2.0 许可
开源了 Engram 实现（\texttt{deepseek-ai/Engram}，
截至 2026 年 3 月获得 3,840 个 star）。
Engram 被认为是
DeepSeek 实现 1M 上下文窗口的关键赋能技术之一。

\begin{keybox}[Engram 与传统方法的对比]
Engram 的条件记忆与已有的长上下文方法是互补的：
\begin{itemize}[noitemsep]
  \item \textbf{vs. 长上下文窗口}：
        上下文窗口适合处理当前会话中的信息，
        Engram 适合存储模型需要“永久记住”的知识
  \item \textbf{vs. RAG}：
        RAG 需要外部检索系统和向量数据库，
        Engram 将检索内置到模型架构中，
        延迟更低且无需外部依赖
  \item \textbf{vs. KV cache 压缩}：
        KV cache 压缩减少了计算中间结果的存储，
        Engram 增加了独立于计算的知识存储
\end{itemize}
\end{keybox}

\subsection{Google Titans：神经记忆架构}

2025 年 12 月，Google 发表了
\textbf{Titans 架构}~\cite{titans2025}，
引入了\textbf{Neural Memory 模块}——
一种可在运行时自我更新的长期记忆机制，
其设计灵感来源于生物大脑的记忆系统。

Titans 的核心架构将模型分为三个组件：
\begin{itemize}[noitemsep]
  \item \textbf{短期注意力核心}（Short-Term Attention Core）：
        标准的 Transformer 注意力层，
        负责处理当前窗口内的精确上下文交互
  \item \textbf{长期神经记忆}（Long-Term Neural Memory）：
        一个在推理时持续更新的神经网络模块，
        能够将跨越数百万 token 的信息
        压缩存储为持久化的记忆表示。
        不同于 RNN 的隐状态，
        神经记忆通过自监督学习目标
        在运行时动态更新其参数——
        真正实现了“边推理边学习”
  \item \textbf{持久记忆}（Persistent Memory）：
        存储不随输入变化的任务知识，
        类似于模型的“先验知识”
\end{itemize}

Titans 的突破性在于：
它能够处理\textbf{数百万 token}的输入，
同时保持跨越数周甚至数月对话的记忆连续性。
Google 将其称为“AI 失忆症的终结”——
传统模型在每次对话结束后
就“忘记”了所有交互内容，
而 Titans 的神经记忆可以在对话间持续积累知识。

从计算效率的角度看，
Titans 结合了 RNN 的处理速度
和 Transformer 的表示精度——
神经记忆模块的更新操作在时间复杂度上
类似于 RNN 的状态更新（$O(1)$ per token），
而短期注意力核心保证了
局部上下文的精确处理。

Titans 与 Engram 代表了长期记忆的两种截然不同的范式。
Engram 是\textbf{静态记忆}——
嵌入表在训练后固定，
推理时只进行读取（查找）操作。
它适合存储“事实性知识”——
“法国的首都是巴黎”这类不随上下文变化的信息。
Titans 是\textbf{动态记忆}——
神经记忆模块在推理时持续更新，
能够学习和存储当前会话中的新信息。
它适合需要“在线学习”的场景——
“用户刚才说他的名字是 Alice，
之后提到 Alice 时要使用这个信息”。

两种范式可以互补：
Engram 提供稳定的、高效的知识存储底层，
Titans 在其上添加动态的、会话级的记忆更新。
这种“静态 + 动态”的双层记忆架构
与人类认知中“长期记忆 + 工作记忆”的分工
有直觉上的对应。

一个开放的研究问题是：
\textbf{动态记忆的更新应该以什么粒度进行？}
Titans 在每个 token 处理后都更新记忆参数，
这提供了最精细的信息捕获能力，
但也带来了“记忆污染”的风险——
模型可能将不重要的信息
或甚至错误的信息写入长期记忆。
一种替代方案是“门控更新”：
只在模型判断当前信息“值得记住”时才更新记忆。
这种选择性更新策略模拟了
人类记忆中“注意力驱动的编码”机制——
我们只记住我们主动关注的信息。

\subsection{Agent 记忆系统的爆发}

随着 LLM Agent（智能体）从研究原型走向生产部署，
Agent 的长期记忆管理在 2026 年初
成为了研究的焦点领域。
与模型层面的长期记忆不同，
Agent 记忆系统关注的是
\textbf{在多次交互和长时间跨度中
管理和利用经验知识}。

\subsubsection{代表性工作}

\textbf{AriadneMem}：
面向“终身 Agent”的记忆架构，
能够将分散在不同时间和不同任务中的
相关证据“穿针引线”般连接起来，
并通过主动的状态更新机制
保持记忆的一致性和时效性。
AriadneMem 的关键设计是一个\textbf{时间关联图}
（temporal association graph）——
记忆不是简单地以列表形式存储，
而是以图结构组织，
其中节点代表记忆片段，
边代表时间关联、因果关联或语义关联。
当 Agent 需要检索记忆时，
它不仅查找与当前 query 最相似的记忆，
还沿着关联边“走”到相关的上下游记忆。
这种图遍历机制使得 Agent 能够
重建被遗忘的推理链——
即使中间步骤的记忆已经模糊，
通过关联边仍可推导出结论。
在多任务持续交互实验中，
AriadneMem 相比平坦记忆存储（flat memory store）
在跨任务知识迁移上提升了约 35\%。

\textbf{PersistBench}：
首个针对 Agent 长期记忆的安全性基准，
揭示了两个重要风险：
（1）\textbf{跨域泄漏}——Agent 在处理任务 A 时
积累的敏感信息可能在任务 B 中被无意泄露；
（2）\textbf{记忆诱导的谄媚}（memory-induced sycophancy）——
Agent 可能过度迎合记忆中记录的用户偏好，
而非提供客观的最优建议。
PersistBench 的测试框架包含 12 个攻击场景，
覆盖了三类威胁模型：
\textbf{被动攻击}（攻击者观察 Agent 的记忆泄露）、
\textbf{主动注入}（攻击者通过精心设计的交互
向 Agent 记忆中植入恶意信息）、
以及\textbf{记忆操纵}（攻击者试图修改或删除 Agent 的已有记忆）。
在 8 个主流 Agent 框架上的评估显示，
没有任何框架能完全抵御所有三类攻击——
这凸显了记忆安全作为一个研究方向的紧迫性。

\textbf{Adaptive Memory Structures}：
根据上下文自适应地选择记忆类型
（工作记忆、情景记忆、语义记忆），
而非使用单一的记忆存储。
这种设计模仿了人类认知心理学中的
多存储记忆模型。
具体来说，\textbf{工作记忆}（working memory）
存储当前任务执行中的中间状态，
容量有限但访问速度最快，
类似于 Transformer 的上下文窗口。
\textbf{情景记忆}（episodic memory）
存储具体交互事件的时间序列——
“用户在 3 月 15 日要求修改了 API 设计”——
保留了时间戳和因果上下文。
\textbf{语义记忆}（semantic memory）
存储从多次交互中抽象出的通用知识——
“这个用户偏好函数式编程风格”——
不绑定到具体事件。
Agent 在每个时间步通过一个
元认知控制器（metacognitive controller）
决定从哪种记忆中检索信息以及将新经验写入哪种记忆。
实验显示，这种多存储架构
比统一记忆存储在 72 小时持续交互测试中
平均减少了 40\% 的冗余记忆写入，
同时将关键信息的召回率提升了 22\%。

\textbf{MemoryLLM}：
将持久化记忆 token 通过自注意力机制
融合到模型的前向传播中。
记忆 token 作为额外的 KV 对
参与每一层的注意力计算，
无需修改模型架构即可实现记忆增强。
MemoryLLM 的实现机制是在每个 Transformer 层的
注意力计算中，将标准的 KV cache
与一组\textbf{持久化记忆键值对}拼接在一起。
这些持久化 KV 对不随上下文窗口滑动而被丢弃——
它们在多次推理调用之间保持不变，
类似于为模型安装了一个“外挂硬盘”。
记忆的写入通过一个门控机制实现：
当模型判断当前输入包含值得长期记住的信息时，
门控信号将该信息编码为新的持久化 KV 对。
这种“即插即用”的设计
使得 MemoryLLM 可以应用于任何标准的
Transformer 模型而不需要重新训练——
只需在推理代码中修改 KV cache 的拼接逻辑即可。
在长篇对话一致性测试中，
MemoryLLM 将跨越 50K token 的事实准确率
从基线的 67\% 提升到 89\%。

2026 年 3 月发表的综述论文
“Memory for Autonomous LLM Agents”
~\cite{agentmemory2026}（arXiv:2603.07670）
对这一快速发展的领域进行了全面梳理，
将 Agent 记忆系统分类为
读取、写入、管理和组织四个维度，
并指出\textbf{记忆安全}将成为
Agent 大规模部署前必须解决的关键问题。

\begin{warnbox}[长期记忆的开放问题]
长期记忆虽然前景广阔，但仍面临多项挑战：
\begin{itemize}[noitemsep]
  \item \textbf{记忆可靠性}：如何确保长期记忆中的信息
        不会随时间“腐化”或产生虚假关联？
  \item \textbf{遗忘机制}：有选择地遗忘过时信息
        与“不可学习”（machine unlearning）的
        隐私合规需求如何协调？
  \item \textbf{评估困难}：如何系统地评估
        跨越数天、数周时间跨度的记忆能力？
        现有的基准测试远未解决这个问题。
        NIAH 类测试只覆盖了“单次会话内的检索”，
        跨会话的记忆持久性和一致性
        缺乏标准化的评估方法
  \item \textbf{安全风险}：PersistBench 已经表明
        长期记忆可能引入新的安全攻击面。
        记忆中毒（memory poisoning）
        可能成为新型攻击向量——
        攻击者通过精心设计的交互
        将错误信息植入 Agent 的长期记忆，
        在未来的交互中误导 Agent 做出错误决策。
        这种攻击特别隐蔽，
        因为它不需要修改模型权重，
        只需污染模型的外部记忆存储
  \item \textbf{隐私合规}：长期记忆如果存储了用户的个人信息，
        如何在用户要求删除时完整清除？
        “被遗忘权”（right to be forgotten）
        在分布式记忆系统中的实现
        是一个未解决的工程和法律问题
\end{itemize}
\end{warnbox}

\begin{corebox}[长上下文的未来方向]
截至 2026 年初，长上下文领域有几个值得关注的研究方向：
\begin{enumerate}[noitemsep]
  \item \textbf{混合架构}：iRoPE（LLaMA 4）和 Transformer-SSM 混合模型
        通过架构创新扩展上下文，
        而非仅依赖位置编码技巧。
        SSM 层不需要 KV cache，
        天然适合超长上下文
  \item \textbf{稀疏注意力训练}：
        DeepSeek 的 NSA 证明了稀疏注意力可以端到端训练，
        不再需要全注意力预训练后蒸馏的两阶段流程。
        Token Sparse Attention、Double-P、DySCO 等后续工作
        在 2026 年初进一步丰富了稀疏注意力的技术工具箱
  \item \textbf{KV cache 压缩}：
        MLA（DeepSeek）、GQA、动态 KV 驱逐等技术
        持续降低长上下文推理的内存开销
  \item \textbf{长期记忆}：
        Engram（DeepSeek）和 Titans（Google）
        开辟了“计算与记忆分离”的新范式，
        Agent 记忆系统的爆发
        进一步将记忆管理推向工程实践
  \item \textbf{无限上下文}：
        Ring Attention + 流水线并行使得
        理论上可以处理任意长度的序列，
        瓶颈从「能否处理」转变为「能否有效利用」
  \item \textbf{Context engineering 工具化}：
        自动上下文管理工具（如智能截断、
        动态检索增强、注意力预算分配）
        可能比扩展上下文窗口本身更有实际价值
\end{enumerate}
\end{corebox}


%% ============================================================
%% NEW REFERENCES (to be added to refs.bib):
%% meta2025llama4 = Meta. "The Llama 4 Herd: The Beginning of a New Era of Natively Multimodal AI Innovation." Meta AI Blog, April 2025.
%% chroma2025contextrot = Hong, Kelly et al. "Context Rot: How Increasing Input Tokens Impacts LLM Performance." Chroma Technical Report, July 2025.
%% liu2023lost = Liu, Nelson et al. "Lost in the Middle: How Language Models Use Long Contexts." arXiv:2307.03172, 2023.
%% anthropic2025contexteng = Anthropic. "Context Engineering for AI Agents." Anthropic Blog, 2025.
%%
%% --- ADDED 2026-03-12 ---
%% engram2026 = DeepSeek-AI. "Conditional Memory via Scalable Lookup: A New Axis of Sparsity for LLMs." arXiv:2601.07372, January 2026.
%% titans2025 = Google Research. "Titans: Learning to Memorize at Test Time." arXiv:2501.00663, December 2025.
%% tokensparse2026 = "Token Sparse Attention: Per-Head Q/K/V Compression for Long-Context Attention." arXiv:2602.03216, February 2026.
%% doublep2026 = "Double-P: Hierarchical Top-p Sparse Attention." arXiv:2602.05191, February 2026.
%% dysco2026 = "DySCO: Dynamic Attention-Scaling via Retrieval Heads." arXiv:2602.22175, February 2026.
%% lostmiddle2026 = "Lost in the Middle: An Exact Theory of Transformer Position Bias." arXiv:2603.10123, March 2026.
%% ultralong2025 = "UltraLong-8B: Efficient Training from 128K to 4M Tokens." arXiv:2504.06214, 2025.
%% agentmemory2026 = "Memory for Autonomous LLM Agents: A Comprehensive Survey." arXiv:2603.07670, March 2026.
%% chen2023extending = Chen et al., "Extending Context Window of Large Language Models via Positional Interpolation", arXiv:2306.15595, 2023
%% llmlingua2024 = Jiang et al., "LLMLingua: Compressing Prompts for Accelerated Inference of Large Language Models", EMNLP 2024
%% xiao2024streamingllm = Xiao et al., "Efficient Streaming Language Models with Attention Sinks", arXiv:2309.17453, 2024
%% autocompressor2023 = Chevalier et al., "AutoCompressor: Learning to Compress Long Contexts", arXiv:2305.14788, 2023
%% mohtashami2023landmark = Mohtashami and Jaggi, "Landmark Attention: Random-Access Infinite Context Length for Transformers", arXiv:2305.16300, 2023
%% hsieh2024ruler = Hsieh et al., "RULER: What's the Real Context Size of Your Language Models?", arXiv:2404.06654, 2024
%% kwon2023vllm = Kwon et al., "Efficient Memory Management for Large Language Model Serving with PagedAttention", SOSP 2023
%%
%% --- ADDED 2026-04-26 (window 2026-03-25 to 2026-04-26) ---
%% deepseekv4 = DeepSeek-AI. "DeepSeek-V4 Technical Report: CSA + HCA + SWA Hybrid Attention for 1M Context." 2026-04-24.
%% deepseekv32dsa = DeepSeek-AI. "DeepSeek-V3.2-Exp: Trainable Sparse Attention with Lightning Indexer." 2025-09-29.
%% nsa2025 = Yuan et al. "Native Sparse Attention." arXiv:2502.11089, ACL 2025.
%% moba2025 = Lu et al. "MoBA: Mixture of Block Attention." arXiv:2502.13189, 2025.
%% nsmla2026 = Anonymous. "NSMLA." OpenReview r9xDZJgIKk, 2026.
%% mamba3hybrid2026 = Together AI. "Mamba-3 hybrid (5 SSM + 1 NoPE Attention)." 2026-03-17.
%% kazemnejad2023nope = Kazemnejad et al. "The Impact of Positional Encoding on Length Generalization." arXiv:2305.19466, 2023.
%% nakanishi2025ssmax = Nakanishi. "Scalable-Softmax." arXiv:2501.19399, 2025.
%% llama4irope2026 = Meta. "Llama 4 iRoPE 256K-to-10M." 2026-04 (with independent evaluations).
%% shuffle2026 = Microsoft. "Shuffle the Context (RoPE-Perturbed Self-Distillation)." arXiv:2604.14339, 2026-04-15.
%% shortdatalongctx2026 = Meta. "Short Data, Long Context." arXiv:2604.06070, 2026-04-07.
%% dashkv2026 = "DASH-KV." arXiv:2604.19351, 2026-04-21.
%% triattention2026 = "TriAttention." arXiv:2604.04921, 2026-04-06.
%% lcakv2026 = "LCA: Latent Context Anchoring." arXiv:2604.12452, 2026.
%% ttkv2026 = "TTKV: Time-Tiered KV Cache." arXiv:2604.19769, 2026-03-27.
%% contextroottheorem2026 = "The Root Theorem of Context Engineering." arXiv:2604.20874, 2026-03-29.
%% opsdl2026 = "OPSDL: On-Policy Self-Distillation for Long Context." arXiv:2604.17535, 2026.
%% structkv2026 = "StructKV." arXiv:2604.06746, 2026.
%% longact2026 = "LongAct." arXiv:2604.14922, 2026.
%% anthropic2026context1m = Anthropic. "1M Context GA." claude.com/blog/1m-context-ga, 2026-03-13.
%% tokenmix2026latency = tokenmix.ai. "1M Prefill Latency Benchmark." 2026-04-20.
%% longbenchv2decomp2026 = "Long Reasoning Decomposition." arXiv:2604.07981, 2026.
%% ============================================================
